@c -*-texinfo-*-

@c  Copyright (c) 1996 Blake McBride
@c  All rights reserved.
@c
@c  Redistribution and use in source and binary forms, with or without
@c  modification, are permitted provided that the following conditions are
@c  met:
@c
@c  1. Redistributions of source code must retain the above copyright
@c  notice, this list of conditions and the following disclaimer.
@c
@c  2. Redistributions in binary form must reproduce the above copyright
@c  notice, this list of conditions and the following disclaimer in the
@c  documentation and/or other materials provided with the distribution.
@c
@c  THIS SOFTWARE IS PROVIDED BY THE COPYRIGHT HOLDERS AND CONTRIBUTORS
@c  "AS IS" AND ANY EXPRESS OR IMPLIED WARRANTIES, INCLUDING, BUT NOT
@c  LIMITED TO, THE IMPLIED WARRANTIES OF MERCHANTABILITY AND FITNESS FOR
@c  A PARTICULAR PURPOSE ARE DISCLAIMED. IN NO EVENT SHALL THE COPYRIGHT
@c  HOLDER OR CONTRIBUTORS BE LIABLE FOR ANY DIRECT, INDIRECT, INCIDENTAL,
@c  SPECIAL, EXEMPLARY, OR CONSEQUENTIAL DAMAGES (INCLUDING, BUT NOT
@c  LIMITED TO, PROCUREMENT OF SUBSTITUTE GOODS OR SERVICES; LOSS OF USE,
@c  DATA, OR PROFITS; OR BUSINESS INTERRUPTION) HOWEVER CAUSED AND ON ANY
@c  THEORY OF LIABILITY, WHETHER IN CONTRACT, STRICT LIABILITY, OR TORT
@c  (INCLUDING NEGLIGENCE OR OTHERWISE) ARISING IN ANY WAY OUT OF THE USE
@c  OF THIS SOFTWARE, EVEN IF ADVISED OF THE POSSIBILITY OF SUCH DAMAGE.

@chapter Library Reference

This chapter gives a detailed description of each WDS class.  It is
organized by facility so that in order to look something up, you would
first go to the section which deals with the facility you are interested
in.  And then the functions are in alphabetical order.

There is a naming convention used with Dynace generics.  All generics
start with either a lower case ``g'', ``v'' or ``m'', and are always
followed by an upper case letter.  The ones which start with ``g'' are
normal generics and may be treated like normal C functions.  The ones
which begin with ``v'' use the variable argument facilities of C and
you should, therefore, take a bit extra care when using them since
there is no compile time argument checking being done with these functions.

Since all generics start with either ``g'', ``v'' or ``m'' and in order to
avoid the difficulty associated with grouping all the generics under
three letters, the first letter is dropped for indexing or heading
purposes.  Therefore if you are looking up a generic, it will always
appear in the index or header with its first letter missing.  The syntax
description and example code, however, will show the entire name.

There is one thing about Dynace which you must be aware of, however.
Dynace generic functions (which is what most of these functions are)
allow the same function to be called in many circumstances.  These
generic functions dynamically dispatch to the appropriate procedures
based on the type of the first argument to the generic function (or
``generic'' for short).  So even though you are calling the same
function name, entirely different functionality may occur based on the
type of the first argument.  This fact is seldom, if ever, a problem
since you know whether you are dealing with an icon or a dialog, for
example, at any given point.  Just be sure to look up the generic in
the appropriate section.


@section A Note To Dynace Language Users
In an effort to hide unnecessary details associated with class vs.
instance methods from WDS users who don't care about the differences,
these differences are not discussed or grouped as they are in the
Dynace manual.  You can always distinguish between the two by what
the method takes as its first argument.  If it's a class, then you
have a class method, otherwise it's an instance method.

If you are a WDS user with interest in the Dynace object oriented
extension to C, it is fully documented in its associated manual.




@section Class Hierarchy
This Dynace Windows Development System contains the following class
hierarchy:

@c @example
@c @group
@iftex
@break

@input wdscls1.tex


@end iftex
@c @end group
@c @end example





@section Application
The @code{Application} class is the class used to control application
wide defaults such as fonts, cursors, brushes and so on.  It is also
the class which gets executed by Windows first and executes your
application specific @code{start} function.

There are no instance methods associated with the @code{Application} class.
All access to this class is through class methods via the global class
object @code{Application}.














@deffn {AddGlobal} AddGlobal::Application
@sp 2
@example
@group
r = gAddGlobal(Application, key, val);

char    *key;   /*  dictionary key     */
object  val;    /*  value to store     */
object  r;      /*  val on success, NULL if key exists  */
@end group
@end example
This class method adds an object to the application-wide global
dictionary under @code{key}.  If @code{key} already exists, the call
fails and returns @code{NULL}; the existing entry is not modified.
@example
@group
@exdent Example:

gAddGlobal(Application, "config", configObj);
@end group
@end example
@sp 1
See also:  @code{GetGlobal, ChangeGlobal, RemoveGlobal, DisposeGlobal}
@end deffn



@deffn {ChangeGlobal} ChangeGlobal::Application
@sp 2
@example
@group
r = gChangeGlobal(Application, key, val);

char    *key;   /*  dictionary key       */
object  val;    /*  new value to store   */
object  r;      /*  result               */
@end group
@end example
This class method replaces the value associated with an existing
@code{key} in the application-wide global dictionary.  The previous
value is not disposed.
@example
@group
@exdent Example:

gChangeGlobal(Application, "config", newConfigObj);
@end group
@end example
@sp 1
See also:  @code{AddGlobal, GetGlobal, RemoveGlobal, DisposeGlobal}
@end deffn



@deffn {CheckMessages} CheckMessages::Application
@sp 2
@example
@group
r = gCheckMessages(Application);

object  r;      /*  Application  */
@end group
@end example
This class method pumps the Windows message queue for timer, cursor,
and paint messages.  This is useful during long-running operations to
keep the UI responsive.
@example
@group
@exdent Example:

gCheckMessages(Application);
@end group
@end example
@end deffn



@deffn {CmdLine} CmdLine::Application
@sp 2
@example
@group
cl = gCmdLine(Application);

char    *cl;    /*  command line   */
@end group
@end example
This method is used to gain access to command line which was used when
the application was launched.
@example
@group
@exdent Example:

int     cl;

cl = gCmdLine(Application);
@end group
@end example
@sp 1
See also:  @code{Instance, Show, PrevInstance}
@end deffn















@c Error::Application removed — gError is a Dynace Object method, not
@c an Application method.  Use gErrorMessage(Application, msg) instead.



@deffn {DisallowMultipleInstances} DisallowMultipleInstances::Application
@sp 2
@example
@group
r = gDisallowMultipleInstances(Application, tag, msg);

char    *tag;   /*  unique application identifier  */
char    *msg;   /*  message to show, or NULL       */
int     r;      /*  non-zero if another instance is running  */
@end group
@end example
This class method checks whether another instance of the application
is already running, using @code{tag} as a system-wide identifier.
If another instance is found and @code{msg} is non-NULL, the message
is displayed.  Returns non-zero if a duplicate is detected.
@example
@group
@exdent Example:

if (gDisallowMultipleInstances(Application, "MyApp", "Already running."))
    return 1;
@end group
@end example
@end deffn



@deffn {DisposeAtExit} DisposeAtExit::Application
@sp 2
@example
@group
r = gDisposeAtExit(Application, obj);

object  obj;    /*  object to dispose at exit  */
object  r;      /*  Application                */
@end group
@end example
This class method registers @code{obj} to be automatically disposed
when the application exits.
@example
@group
@exdent Example:

gDisposeAtExit(Application, myResource);
@end group
@end example
@sp 1
See also:  @code{RemoveAtExit}
@end deffn



@deffn {DisposeGlobal} DisposeGlobal::Application
@sp 2
@example
@group
r = gDisposeGlobal(Application, key);

char    *key;   /*  dictionary key  */
object  r;      /*  Application on success, NULL if not found  */
@end group
@end example
This class method removes an entry from the application-wide global
dictionary and deep-disposes both the key and the value object.
@example
@group
@exdent Example:

gDisposeGlobal(Application, "config");
@end group
@end example
@sp 1
See also:  @code{AddGlobal, GetGlobal, ChangeGlobal, RemoveGlobal}
@end deffn



@deffn {ErrorMessage} ErrorMessage::Application
@sp 2
@example
@group
r = gErrorMessage(Application, msg);

char    *msg;   /*  error message text  */
object  r;      /*  Application         */
@end group
@end example
This class method displays @code{msg} in a modal error message box.
@example
@group
@exdent Example:

gErrorMessage(Application, "File not found.");
@end group
@end example
@sp 1
See also:  @code{Message, MessageMode}
@end deffn



@deffn {GetArg} GetArg::Application
@sp 2
@example
@group
s = gGetArg(Application, num);

int     num;    /*  argument index (0-based)  */
char    *s;     /*  argument string           */
@end group
@end example
This class method returns the command-line argument at index
@code{num} (0-based).  Returns an empty string if @code{num} is out
of range.
@example
@group
@exdent Example:

char *prog = gGetArg(Application, 0);
@end group
@end example
@sp 1
See also:  @code{NumArgs}
@end deffn



@deffn {GetBackBrush} GetBackBrush::Application
@sp 2
@example
@group
bo = gGetBackBrush(Application);

object  bo;     /*  brush object  */
@end group
@end example
This method is used to obtain a copy of the default background brush
object associated with the application.  The returned object will be an
instance of one of the subclasses of the @code{Brush} class.

The object returned must, either explicitly or implicitly, be
disposed when it is no longer needed.  This is normally done
automatically by WDS when it is associated with a window.

See the @code{Brush} class and its subclasses for further details.
@example
@group
@exdent Example:

object  bo;

bo = gGetBackBrush(Application);
@end group
@end example
@sp 1
See also:  @code{GetTextBrush, SetBackBrush}
@end deffn



@deffn {GetColorTheme} GetColorTheme::Application
@sp 2
@example
@group
r = gGetColorTheme(Application, rp, gp, bp);

int     *rp;    /*  pointer to receive red value    */
int     *gp;    /*  pointer to receive green value  */
int     *bp;    /*  pointer to receive blue value   */
int     r;      /*  non-zero if a theme is active   */
@end group
@end example
This class method retrieves the current color theme RGB values.
Returns non-zero if a color theme has been set.
@example
@group
@exdent Example:

int r, g, b;
if (gGetColorTheme(Application, &r, &g, &b))
    /* theme is active */;
@end group
@end example
@sp 1
See also:  @code{SetColorTheme}
@end deffn



@deffn {GetCursor} GetCursor::Application
@sp 2
@example
@group
co = gGetCursor(Application);

object  co;     /*  cursor object  */
@end group
@end example
This method is used to obtain a copy of the default cursor object
associated with the application.  The returned object will be an
instance of one of the subclasses of the @code{Cursor} class.

The object returned must, either explicitly or implicitly, be
disposed when it is no longer needed.  This is normally done
automatically by WDS when it is associated with a window.

See the @code{Cursor} class and its subclasses for further details.
@example
@group
@exdent Example:

object  co;

co = gGetCursor(Application);
@end group
@end example
@sp 1
See also:  @code{SetCursor}
@end deffn



@deffn {GetFont} GetFont::Application
@sp 2
@example
@group
fo = gGetFont(Application);

object  fo;     /*  font object  */
@end group
@end example
This method is used to obtain a copy of the default font object
associated with the application.  The returned object will be an
instance of one of the subclasses of the @code{Font} class.

The object returned must, either explicitly or implicitly, be
disposed when it is no longer needed.  This is normally done
automatically by WDS when it is associated with a window.

See the @code{Font} class and its subclasses for further details.
@example
@group
@exdent Example:

object  fo;

fo = gGetFont(Application);
@end group
@end example
@sp 1
See also:  @code{SetFont}
@end deffn



@deffn {GetGlobal} GetGlobal::Application
@sp 2
@example
@group
r = gGetGlobal(Application, key);

char    *key;   /*  dictionary key  */
object  r;      /*  stored object, or NULL  */
@end group
@end example
This class method retrieves an object previously stored in the
application-wide global dictionary under @code{key}.  Returns
@code{NULL} if @code{key} is not found.
@example
@group
@exdent Example:

object cfg = gGetGlobal(Application, "config");
@end group
@end example
@sp 1
See also:  @code{AddGlobal, ChangeGlobal, RemoveGlobal, DisposeGlobal}
@end deffn



@deffn {GetIcon} GetIcon::Application
@sp 2
@example
@group
io = gGetIcon(Application);

object  io;     /*  icon object  */
@end group
@end example
This method is used to obtain a copy of the default icon object
associated with the application.  The returned object will be an
instance of one of the subclasses of the @code{Icon} class.

The object returned must, either explicitly or implicitly, be
disposed when it is no longer needed.  This is normally done
automatically by WDS when it is associated with a window.

See the @code{Icon} class and its subclasses for further details.
@example
@group
@exdent Example:

object  io;

io = gGetIcon(Application);
@end group
@end example
@sp 1
See also:  @code{SetIcon}
@end deffn



@deffn {GetName} GetName::Application
@sp 2
@example
@group
nm = gGetName(Application);

char    *nm;    /*  application name  */
@end group
@end example
This method is used to obtain the global name associated with the application.
@example
@group
@exdent Example:

char    *nm;

nm = gGetName(Application);
@end group
@end example
@sp 1
See also:  @code{SetName}
@end deffn



@deffn {GetScalingMode} GetScalingMode::Application
@sp 2
@example
@group
m = gGetScalingMode(Application);

int     m;      /*  scaling mode   */
@end group
@end example
This method is used to get the current scaling mode in effect.
See @code{SetScalingMode} for further details.
@example
@group
@exdent Example:

int     mode;

mode = gGetScalingMode(Application);
@end group
@end example
@sp 1
See also:  @code{SetScalingMode, ScaleToPixels, ScaleToCurrentMode}
@end deffn



@deffn {GetSize} GetSize::Application
@sp 2
@example
@group
r = gGetSize(Application, vert, horz);

int    *vert;   /*  vertical size    */
int    *horz;   /*  horizontal size  */
object  r;      /*  Application      */
@end group
@end example
This method is used to get the total size of the user's screen.
@code{vert} and @code{horz} are in increments dictated by the
mode selected by @code{SetScalingMode}.  
@example
@group
@exdent Example:

int     y, x;

gGetSize(Application, &y, &x);
@end group
@end example
@sp 1
See also:  @code{SetScalingMode, GetSize::Window}
@end deffn



@deffn {GetTextBrush} GetTextBrush::Application
@sp 2
@example
@group
bo = gGetTextBrush(Application);

object  bo;     /*  brush object  */
@end group
@end example
This method is used to obtain a copy of the default text brush object
associated with the application.  The returned object will be an
instance of one of the subclasses of the @code{Brush} class.

The object returned must, either explicitly or implicitly, be
disposed when it is no longer needed.  This is normally done
automatically by WDS when it is associated with a window.

See the @code{Brush} class and its subclasses for further details.
@example
@group
@exdent Example:

object  bo;

bo = gGetTextBrush(Application);
@end group
@end example
@sp 1
See also:  @code{GetBackBrush, SetTextBrush}
@end deffn



@deffn {Instance} Instance::Application
@sp 2
@example
@group
ins = gInstance(Application);

HINSTANCE  ins;    /*  instance handle   */
@end group
@end example
This method is used to gain access to the Windows instance handle
associated with the application.  This handle is mainly used internally
by Windows and WDS, and should not normally be needed by a WDS
programmer.
@example
@group
@exdent Example:

HINSTANCE  h;

h = gInstance(Application);
@end group
@end example
@sp 1
See also:  @code{PrevInstance, CmdLine, Show}
@end deffn



@deffn {IsExiting} IsExiting::Application
@sp 2
@example
@group
r = gIsExiting(Application);

int     r;      /*  1 if exiting, 0 otherwise  */
@end group
@end example
This class method returns 1 if the application is in the process of
exiting, 0 otherwise.
@example
@group
@exdent Example:

if (gIsExiting(Application))
    return;
@end group
@end example
@end deffn



@deffn {Language} Language::Application
@sp 2
@example
@group
lang = gLanguage(Application);

int     lang;   /*  current language index  */
@end group
@end example
This class method returns the current language index previously set
by @code{SetLanguage}.  The default is 0 (English).
@example
@group
@exdent Example:

int lang = gLanguage(Application);
@end group
@end example
@sp 1
See also:  @code{SetLanguage}
@end deffn



@deffn {Message} Message::Application
@sp 2
@example
@group
r = gMessage(Application, msg);

char    *msg;   /*  message text  */
object  r;      /*  Application   */
@end group
@end example
This class method displays @code{msg} in a message box, provided
message mode is enabled.  If message mode is disabled, no dialog
appears.
@example
@group
@exdent Example:

gMessage(Application, "Operation complete.");
@end group
@end example
@sp 1
See also:  @code{ErrorMessage, MessageMode}
@end deffn



@deffn {MessageMode} MessageMode::Application
@sp 2
@example
@group
prev = gMessageMode(Application, mode);

int     mode;   /*  >= 0 to set, < 0 to query  */
int     prev;   /*  previous mode value         */
@end group
@end example
This class method gets or sets whether @code{gMessage} shows a
message box.  Pass a non-negative value to set the mode, or a
negative value to query without changing it.  Returns the previous
mode.
@example
@group
@exdent Example:

int old = gMessageMode(Application, 0);  /* disable messages */
@end group
@end example
@sp 1
See also:  @code{Message, ErrorMessage}
@end deffn



@deffn {NumArgs} NumArgs::Application
@sp 2
@example
@group
n = gNumArgs(Application);

int     n;      /*  number of command-line arguments  */
@end group
@end example
This class method returns the number of command-line arguments
passed to the application.
@example
@group
@exdent Example:

int n = gNumArgs(Application);
@end group
@end example
@sp 1
See also:  @code{GetArg}
@end deffn



@deffn {OpenURL} OpenURL::Application
@sp 2
@example
@group
r = gOpenURL(Application, url);

char    *url;   /*  URL or file path  */
object  r;      /*  Application       */
@end group
@end example
This class method opens a URL or file using the system default
handler via @code{ShellExecute}.  Control returns immediately.
@example
@group
@exdent Example:

gOpenURL(Application, "https://example.com");
@end group
@end example
@sp 1
See also:  @code{OpenURLWait}
@end deffn



@deffn {OpenURLWait} OpenURLWait::Application
@sp 2
@example
@group
r = gOpenURLWait(Application, url);

char    *url;   /*  URL or file path  */
object  r;      /*  Application       */
@end group
@end example
This class method opens a URL or file using the system default
handler and blocks until the spawned process finishes.
@example
@group
@exdent Example:

gOpenURLWait(Application, "report.pdf");
@end group
@end example
@sp 1
See also:  @code{OpenURL}
@end deffn



@deffn {PrevInstance} PrevInstance::Application
@sp 2
@example
@group
ins = gPrevInstance(Application);

HINSTANCE  ins;    /*  instance handle   */
@end group
@end example
This method is used to gain access to the Windows instance handle
associated with the previous instance of this application, should more
than one be running.  This handle is mainly used internally by Windows
and WDS, and should not normally be needed by a WDS programmer.

This value will always be @code{NULL} under Windows NT.
@example
@group
@exdent Example:

HINSTANCE  h;

h = gPrevInstance(Application);
@end group
@end example
@sp 1
See also:  @code{Instance, CmdLine, Show}
@end deffn



@deffn {QuitApplication} QuitApplication::Application
@sp 2
@example
@group
r = gQuitApplication(Application, ret);

int     ret;    /*  app return value  */
object  r;      /*  Application       */
@end group
@end example
This method is used to terminate an application.  The value passed will
be used as the return value of the application.
@example
@group
@exdent Example:

gQuitApplication(Application, 0);
@end group
@end example
@sp 1
See also:  @code{Error}
@end deffn



@deffn {RemoveAtExit} RemoveAtExit::Application
@sp 2
@example
@group
r = gRemoveAtExit(Application, obj);

object  obj;    /*  object to remove from exit set  */
object  r;      /*  Application                     */
@end group
@end example
This class method removes @code{obj} from the set of objects that
will be disposed at application exit.  The object itself is not
disposed.
@example
@group
@exdent Example:

gRemoveAtExit(Application, myResource);
@end group
@end example
@sp 1
See also:  @code{DisposeAtExit}
@end deffn



@deffn {RemoveGlobal} RemoveGlobal::Application
@sp 2
@example
@group
r = gRemoveGlobal(Application, key);

char    *key;   /*  dictionary key  */
object  r;      /*  Application on success, NULL if not found  */
@end group
@end example
This class method removes an entry from the application-wide global
dictionary without disposing the value object.  Use this when you
want to retain the value after removal.
@example
@group
@exdent Example:

gRemoveGlobal(Application, "config");
@end group
@end example
@sp 1
See also:  @code{AddGlobal, GetGlobal, ChangeGlobal, DisposeGlobal}
@end deffn



@deffn {ReturnAsTab} ReturnAsTab::Application
@sp 2
@example
@group
prev = gReturnAsTab(Application, mode);

int     mode;   /*  >= 0 to set, < 0 to query  */
int     prev;   /*  previous mode value         */
@end group
@end example
This class method gets or sets whether the Return key acts like the
Tab key for dialog navigation.  Pass a non-negative value to set
the mode, or a negative value to query without changing it.  Returns
the previous mode.
@example
@group
@exdent Example:

gReturnAsTab(Application, 1);  /* enable Return-as-Tab */
@end group
@end example
@end deffn



@deffn {ScaleToCurrentMode} ScaleToCurrentMode::Application
@sp 2
@example
@group
r = gScaleToCurrentMode(Application, y, x, fnt);

int     *y;     /*  row position         */
int     *x;     /*  column position      */
object  fnt;    /*  current font object  */
object  r;      /*  Application          */
@end group
@end example
This method is used to convert coordinates from pixels to the current
scaling mode.  It is mainly used internally by WDS in order to
convert Window's standard pixel positions into the current scaling mode.

On entry @code{x} and @code{y} point to values which are coordinates in
terms of pixels.  After this method returns, their value will be changed
to be in terms of the current scaling mode (see @code{SetScalingMode}).

@code{fnt} is a font object which is used only if the current scaling
mode is relative to a font.  If so, the font it will be related to
will be the one passed.
@example
@group
@exdent Example:

int     y, x;
object  fnt;

y = some position;
x = some position;
fnt = some font object;
gScaleToCurrentMode(Application, &y, &x, fnt);
@end group
@end example
@sp 1
See also:  @code{SetScalingMode, GetScalingMode, ScaleToPixels}
@end deffn



@deffn {ScaleToPixels} ScaleToPixels::Application
@sp 2
@example
@group
r = gScaleToPixels(Application, y, x, fnt);

int     *y;     /*  row position         */
int     *x;     /*  column position      */
object  fnt;    /*  current font object  */
object  r;      /*  Application          */
@end group
@end example
This method is used to convert coordinates from the current scaling mode
to pixels.  It is mainly used internally by WDS in order to convert
your coordinates into standard pixel positions.

On entry @code{x} and @code{y} point to values which are coordinates
in terms of the current scaling mode (see @code{SetScalingMode}).
After this method returns, their value will be changed to be in terms
of pixels.

@code{fnt} is a font object which is used only if the current scaling
mode is relative to a font.  If so, the font it will be related to
will be the one passed.
@example
@group
@exdent Example:

int     y, x;
object  fnt;

y = some position;
x = some position;
fnt = some font object;
gScaleToPixels(Application, &y, &x, fnt);
@end group
@end example
@sp 1
See also:  @code{SetScalingMode, GetScalingMode, ScaleToCurrentMode}
@end deffn



@deffn {SetBackBrush} SetBackBrush::Application
@sp 2
@example
@group
r = gSetBackBrush(Application, bo);

object  bo;     /*  brush object         */
object  r;      /*  brush object passed  */
@end group
@end example
This method is used to set the application wide default background
brush.  This is the brush used to color everything except the text that
appears on windows and dialogs.  @code{bo} must be an instance of one of
the subclasses of @code{Brush}.  All windows and dialogs will use the
application wide default brush which is in effect when they are created
unless a specific brush object is specified for a particular window.

When a new default brush object is set, any previous default object
will be disposed.  If no default is set, WDS uses the system brush
identified as @code{COLOR_WINDOW}.

See the @code{Brush} class and its subclasses for further details.
@example
@group
@exdent Example:

gSetBackBrush(Application, vNew(SystemBrush, COLOR_WINDOW));
@end group
@end example
@sp 1
See also:  @code{GetBackBrush, SetTextBrush}
@end deffn



@deffn {SetColorTheme} SetColorTheme::Application
@sp 2
@example
@group
gSetColorTheme(Application, r, g, b);

int     r;      /*  red component (0-255)    */
int     g;      /*  green component (0-255)  */
int     b;      /*  blue component (0-255)   */
@end group
@end example
This class method sets the application background color theme using
RGB values.
@example
@group
@exdent Example:

gSetColorTheme(Application, 220, 220, 255);
@end group
@end example
@sp 1
See also:  @code{GetColorTheme}
@end deffn



@deffn {SetCursor} SetCursor::Application
@sp 2
@example
@group
r = gSetCursor(Application, co);

object  co;     /*  cursor object         */
object  r;      /*  cursor object passed  */
@end group
@end example
This method is used to set the application wide default cursor.
@code{co} must be an instance of one of the subclasses of @code{Cursor}.
All windows will use the application wide default cursor which is in
effect when they are created unless a specific cursor object is
specified for a particular window.

When a new default cursor object is set, any previous default object
will be disposed.  If no default is set, WDS uses the system cursor
identified as @code{IDC_ARROW}.

See the @code{Cursor} class and its subclasses for further details.
@example
@group
@exdent Example:

gSetCursor(Application, gLoadSys(SystemCursor, IDC_ARROW));
@end group
@end example
@sp 1
See also:  @code{GetCursor, LoadSys::SystemCursor, Load::ExternalCursor}
@end deffn



@deffn {SetFont} SetFont::Application
@sp 2
@example
@group
r = gSetFont(Application, fo);

object  fo;     /*  font object         */
object  r;      /*  font object passed  */
@end group
@end example
This method is used to set the application wide default font.
@code{fo} must be an instance of one of the subclasses of @code{Font}.
All windows will use the application wide default font which is in
effect when they are created unless a specific font object is
specified for a particular window.

When a new default font object is set, any previous default object
will be disposed.  If no default is set, WDS uses the system font
identified as @code{SYSTEM_FONT}.

See the @code{Font} class and its subclasses for further details.
@example
@group
@exdent Example:

gSetFont(Application, vNew(SystemFont, SYSTEM_FONT));
@end group
@end example
@sp 1
See also:  @code{GetFont, Load::SystemFont, New::ExternalFont}
@end deffn



@deffn {SetIcon} SetIcon::Application
@sp 2
@example
@group
r = gSetIcon(Application, io);

object  io;     /*  icon object         */
object  r;      /*  icon object passed  */
@end group
@end example
This method is used to set the application wide default icon.  An icon
associated to a window is the one which is displayed when the window is
iconized.  @code{io} must be an instance of one of the subclasses of
@code{Icon}.  All windows will use the application wide default icon
which is in effect when they are created unless a specific icon object
is specified for a particular window.

When a new default icon object is set, any previous default object
will be disposed.  If no default is set, WDS uses the system icon
identified as @code{IDI_APPLICATION}.

See the @code{Icon} class and its subclasses for further details.
@example
@group
@exdent Example:

gSetIcon(Application, gLoadSys(SystemIcon, IDI_APPLICATION));
@end group
@end example
@sp 1
See also:  @code{GetIcon, LoadSys::SystemIcon, Load::ExternalIcon}
@end deffn



@deffn {SetLanguage} SetLanguage::Application
@sp 2
@example
@group
r = gSetLanguage(Application, lang);

int     lang;   /*  language index (0 = English)  */
object  r;      /*  Application                   */
@end group
@end example
This class method sets the current application language.  Language
index 0 is English; other indices are application-defined.
@example
@group
@exdent Example:

gSetLanguage(Application, 1);  /* Spanish */
@end group
@end example
@sp 1
See also:  @code{Language}
@end deffn



@deffn {SetName} SetName::Application
@sp 2
@example
@group
gSetName(Application, nm);

char    *nm;    /*  application name  */
@end group
@end example
This method is used to give the application a globally accessible name.
This name is accessible via @code{GetName}.  There is no other use
made of this information.
@example
@group
@exdent Example:

gSetName(Application, "My App");
@end group
@end example
@sp 1
See also:  @code{GetName}
@end deffn



@deffn {SetScalingMode} SetScalingMode::Application
@sp 2
@example
@group
r = gSetScalingMode(Application, m);

int     m;      /*  scaling mode   */
int     r;      /*  previous mode  */
@end group
@end example
This method is used to set the scaling mode used by all other WDS methods
which take coordinate positions.  Valid modes are as follows:
@table @code
@item SM_1_PER_CHAR
This mode causes each position to be in increments determined by the
size of the current font.  For example row 7 would mean 7 times the
height of the current font.  Similar to line positions.
@item SM_10_PER_SYSCHAR
This mode causes each position to be in increments determined by one
tenth the size of the system font declared globally to the user's
Windows environment.  For example row 70 would mean 7 times the height
of the Windows global system font.  This allows positioning relative
to a scaling factor determined by the user.
@item SM_PIXELS
This mode performs no conversion.  The application is able to use
pixel coordinated directly.
@end table

The default value is @code{SM_1_PER_CHAR}.

The value returned is the mode which was previously set.
@example
@group
@exdent Example:

gSetScalingMode(Application, SM_PIXELS);
@end group
@end example
@sp 1
See also:  @code{GetScalingMode, ScaleToPixels, ScaleToCurrentMode}
@end deffn



@deffn {SetTextBrush} SetTextBrush::Application
@sp 2
@example
@group
r = gSetTextBrush(Application, bo);

object  bo;     /*  brush object         */
object  r;      /*  brush object passed  */
@end group
@end example
This method is used to set the application wide default text brush.
This is the brush used by all text output (or foreground) to windows or
dialogs.  @code{bo} must be an instance of one of the subclasses of
@code{Brush}.  All windows and dialogs will use the application wide
default brush which is in effect when they are created unless a specific
brush object is specified for a particular window.

When a new default brush object is set, any previous default object
will be disposed.  If no default is set, WDS uses the system brush
identified as @code{COLOR_WINDOWTEXT}.

See the @code{Brush} class and its subclasses for further details.
@example
@group
@exdent Example:

gSetTextBrush(Application, vNew(SystemBrush, COLOR_WINDOWTEXT));
@end group
@end example
@sp 1
See also:  @code{GetTextBrush, SetBackBrush}
@end deffn



@deffn {Show} Show::Application
@sp 2
@example
@group
sv = gShow(Application);

int     sv;     /*  show value   */
@end group
@end example
This method is used to gain access to the show value supplied by Windows
when an application starts.  It is normally used to determine how an
initial application's main window should be shown.  The available
options are Windows macros which begin with @code{SW_} and are fully
documented by the Windows documentation under the @code{WinMain}
function.
@example
@group
@exdent Example:

int     sv;

sv = gShow(Application);
@end group
@end example
@sp 1
See also:  @code{Instance, CmdLine, PrevInstance}
@end deffn



@section Windows
This section documents the @code{Window} class which contains all the
functionality which is common to Main, Child, and Popup windows.  See
chapter 2 of this manual for a description of the different window
types.






@deffn {AddCompletionFunctionAfter} AddCompletionFunctionAfter::Window
@sp 2
@example
@group
r = gAddCompletionFunctionAfter(wind, fun);

object  wind;      /*  a window object     */
int    (*fun)();   /*  completion function  */
object  r;         /*  wind                */
@end group
@end example
This method adds a completion function @code{fun} to the end of the
list of completion functions associated with window @code{wind}.
The completion function is called when the window is closed.
@code{fun} receives the window object as its argument and should
return 0 to cancel the close or non-zero to allow it.

The window object passed is returned.
@example
@group
@exdent Example:

gAddCompletionFunctionAfter(myWind, my_cleanup);
@end group
@end example
@sp 1
See also:  @code{CompletionFunction, AddCompletionFunctionBefore}
@end deffn




@deffn {AddCompletionFunctionBefore} AddCompletionFunctionBefore::Window
@sp 2
@example
@group
r = gAddCompletionFunctionBefore(wind, fun);

object  wind;      /*  a window object     */
int    (*fun)();   /*  completion function  */
object  r;         /*  wind                */
@end group
@end example
This method adds a completion function @code{fun} to the beginning of
the list of completion functions associated with window @code{wind}.
See @code{AddCompletionFunctionAfter} for full details.

The window object passed is returned.
@sp 1
See also:  @code{CompletionFunction, AddCompletionFunctionAfter}
@end deffn




@deffn {AddHandlerAfter} AddHandlerAfter::Window
@sp 2
@example
@group
r = gAddHandlerAfter(wind, msg, func);

object   wind;     /*  a window object  */
unsigned msg;      /*  message          */
long    (*func)(); /*  function pointer */
object  r;         /*  the window obj   */
@end group
@end example
This method is used to associate function @code{func} with Windows window
message @code{msg} for window @code{wind}.  Whenever window @code{wind}
receives message @code{msg}, @code{func} will be called.

@code{wind} is the window object who's messages you wish to process.
@code{msg} is the particular message you wish to trap.  These messages
are fully documented in the Windows documentation in the Messages
section.  They normally begin with @code{WM_}.

@code{func} is the function which gets called whenever the specified
message gets received and takes the following form:
@example
@group
long    func(object     wind,
             HWND       hwnd, 
             UINT       mMsg, 
             WPARAM     wParam, 
             LPARAM     lParam)
@{
        .
        .
        .
        return 0L;  /* or whatever is appropriate  */
@}
@end group
@end example
Where @code{wind} is the window being sent the message.  The remaining
arguments and return value is fully documented in the Windows documentation
under the @code{WindowProc} function and the Windows Messages documentation.

WDS keeps a list of functions associated with each message associated
with each window.  When a particular message is received the appropriate
list of handler functions gets executed sequentially.
@code{AddHandlerAfter} appends the new function to the end of this list,
and @code{AddHandlerBefore} adds the new function to the beginning of
the list.

WDS may also, and optionally, execute the Windows default procedure
associated with a given message either before or after the user added
list of functions.  This behavior may be controlled via
@code{DefaultProcessingMode}.

Windows will only see the return value of the last message handler executed
including, if applicable, the default.
@example
@group
@exdent Example:

int     hSize, vSize;

static  long    process_wm_size(object  wind, 
                                HWND    hwnd, 
                                UINT    mMsg, 
                                WPARAM  wParam, 
                                LPARAM  lParam)
@{
        hSize = LOWORD(lParam);
        vSize = HIWORD(lParam);
        return 0L;
@}

        .
        .
        gAddHandlerAfter(wind, (unsigned) WM_SIZE, process_wm_size);
        .
        .
@end group
@end example
@sp 1
See also:  @code{DefaultProcessingMode, AddHandlerBefore}
@end deffn






@deffn {AddHandlerBefore} AddHandlerBefore::Window
@sp 2
@example
@group
r = gAddHandlerBefore(wind, msg, func);

object   wind;     /*  a window object  */
unsigned msg;      /*  message          */
long    (*func)(); /*  function pointer */
object  r;         /*  the window obj   */
@end group
@end example
This function is fully documented under @code{AddHandlerAfter}.
@c @example
@c @group
@c @exdent Example:
@c @end group
@c @end example
@sp 1
See also:  @code{AddHandlerAfter}
@end deffn







@deffn {AddInitFunctionAfter} AddInitFunctionAfter::Window
@sp 2
@example
@group
r = gAddInitFunctionAfter(wind, fun);

object  wind;      /*  a window object       */
int    (*fun)();   /*  initialization function */
object  r;         /*  wind                  */
@end group
@end example
This method adds an initialization function @code{fun} to the end of
the list of init functions associated with window @code{wind}.
Init functions are called when the window is initialized.
@code{fun} receives the window object as its argument.

The window object passed is returned.
@sp 1
See also:  @code{InitFunction, AddInitFunctionBefore}
@end deffn




@deffn {AddInitFunctionBefore} AddInitFunctionBefore::Window
@sp 2
@example
@group
r = gAddInitFunctionBefore(wind, fun);

object  wind;      /*  a window object       */
int    (*fun)();   /*  initialization function */
object  r;         /*  wind                  */
@end group
@end example
This method adds an initialization function @code{fun} to the beginning of
the list of init functions associated with window @code{wind}.
See @code{AddInitFunctionAfter} for full details.

The window object passed is returned.
@sp 1
See also:  @code{InitFunction, AddInitFunctionAfter}
@end deffn




@deffn {AddTimer} AddTimer::Window
@sp 2
@example
@group
r = gAddTimer(wind, interval, fun);

object  wind;      /*  a window object         */
int     interval;  /*  timer interval in ms    */
int    (*fun)();   /*  timer callback function */
int     r;         /*  timer ID                */
@end group
@end example
This method adds a recurring timer to window @code{wind}.
The function @code{fun} will be called every @code{interval}
milliseconds.  The value returned is a timer ID which can be
passed to @code{KillTimer} to remove the timer.
@example
@group
@exdent Example:

int  timerID;

timerID = gAddTimer(myWind, 1000, my_timer_func);
@end group
@end example
@sp 1
See also:  @code{KillTimer}
@end deffn




@deffn {Associate} Associate::Window
@sp 2
@example
@group
r = mAssociate(wind, itm, fun);

object  wind;   /*  a window object     */
int     itm;    /*  menu item           */
long    (*fun)();  /*  function         */
object   r;        /*  the menu         */
@end group
@end example
This method is used to associate an application specific function
(@code{fun}) with a menu item (@code{itm}) associated with the current
menu attached to window @code{wind}.  Once this is done, if the user selects
the menu option identified by @code{itm}, then the function @code{fun}
will be executed.

@code{itm} is a programmer defined macro which identifies one particular
choice among those available within the menu which is currently attached
to window @code{wind}.  This macro is defined while the programmer defines
the entire menu using the resource editor.

@code{fun} is the function which will be executed when the user selects
menu option @code{itm} and has the following form:
@example
@group
long    fun(object wind, unsigned id)
@{
        .
        .
        .
        return 0L;
@}
@end group
@end example
The function executed (@code{fun}) is passed the window object and
the specific menu id which the user selected and
returns a long.  The return value is documented in the Windows documentation
under the message named @code{WM_COMMAND}.  It should normally be @code{0L}.
@example
@group
@exdent Example:

static  long    file_message(object wind)
@{
        gMessage(wind, "File_Message");
        return 0L;
@}

        .
        .
        mLoadMenu(win, IDR_MENU1);
        mAssociate(win, ID_FILE_MESSAGE, file_message);
@end group
@end example
@sp 1
See also:  @code{LoadMenu, MenuItemMode}
@end deffn











@deffn {AutoDispose} AutoDispose::Window
@sp 2
@example
@group
r = gAutoDispose(wind, mode);

object   wind;     /*  a window object    */
int      mode;     /*  auto dispose mode  */
int      r;        /*  previous value     */
@end group
@end example
This method is used to enable or disable the auto dispose mechanism
associated with a window.  This feature, when enabled, causes WDS to
automatically dispose of the WDS object associated with a window
whenever the user closes the window (@code{wind}).  Normally, when
this feature is disabled, if the user closes a window, WDS removes
the window from the display, but the WDS object remains intact until
the program manually disposes of the window object (via @code{Dispose}).
This feature is most commonly used with asynchronous, popup windows.

Note that it is invalid to attempt to use a WDS object subsequent to it
being disposed, and each unneeded WDS object must be deleted.

@code{mode} is set to @code{1} to enable the feature and @code{0}
to disable it.  The value returned is the prior mode.
@example
@group
@exdent Example:

gAutoDispose(wind, 1);
@end group
@end example
@sp 1
See also:  @code{Dispose}
@end deffn











@deffn {AutoShow} AutoShow::Window
@sp 2
@example
@group
r = gAutoShow(wind, mode);

object   wind;     /*  a window object    */
int      mode;     /*  auto show mode     */
int      r;        /*  previous value     */
@end group
@end example
This method is used to enable or disable the auto show mechanism
associated with a window.  This feature, when enabled, causes WDS to
automatically display (execute @code{Show}) on a window which is written
to.  This enables the programmer to create a window which will automatically
popup the first time the program displays text to it.

@code{mode} is set to @code{1} to enable the feature and @code{0}
to disable it.  The value returned is the prior mode.
@example
@group
@exdent Example:

gAutoShow(wind, 1);
@end group
@end example
@sp 1
See also:  @code{Show}
@end deffn











@deffn {BackBrush} BackBrush::Window
@sp 2
@example
@group
r = gBackBrush(wind, brsh);

object  wind;   /*  a window object   */
object  brsh;   /*  brush object      */
object  r;      /*  wind              */
@end group
@end example
This method is used to determine what brush object is used for the
background of window @code{wind} and performs the same function as the
@code{Use} method when used with a @code{Brush} object.  Any previously
associated brush object will be disposed.  This brush object will also
be automatically disposed when the window is disposed.

The window passed is returned.
@example
@group
@exdent Example:

object  myWind;

gBackBrush(myWind, vNew(SolidBrush, 0, 0, 255));
@end group
@end example
@sp 1
See also:  @code{TextBrush, Use, SetTextBrush::Application}
        and the @code{Brush} classes
@end deffn












@deffn {Center} Center::Window
@sp 2
@example
@group
r = gCenter(wind);

object  wind;   /*  a window object  */
object  r;      /*  wind             */
@end group
@end example
This method centers window @code{wind} on the screen.  If the
window is already shown, it is centered immediately; otherwise
it will be centered when shown.

The window object passed is returned.
@example
@group
@exdent Example:

object  myWind;

myWind = vNew(MainWindow, "My App");
gCenter(myWind);
@end group
@end example
@c @sp 1
@c See also:  @code{}
@end deffn




@deffn {CompletionFunction} CompletionFunction::Window
@sp 2
@example
@group
r = gCompletionFunction(wind, fun);

object  wind;      /*  a window object     */
int    (*fun)();   /*  completion function  */
object  r;         /*  wind                */
@end group
@end example
This method sets the completion function associated with window
@code{wind} to @code{fun}.  The completion function is called
when the window is closed.  @code{fun} receives the window
object as its argument and should return 0 to cancel the
close or non-zero to allow it.

Any previously set completion function is replaced.

The window object passed is returned.
@example
@group
@exdent Example:

static  int  my_completion(object wind)
@{
        /* perform cleanup */
        return 1;  /* allow close */
@}

        .
        .
        gCompletionFunction(myWind, my_completion);
@end group
@end example
@sp 1
See also:  @code{AddCompletionFunctionAfter, AddCompletionFunctionBefore}
@end deffn




@deffn {DefaultProcessingMode} DefaultProcessingMode::Window
@sp 2
@example
@group
r = gDefaultProcessingMode(wind, msg, mode);

object   wind;     /*  a window object          */
unsigned msg;      /*  message                  */
int      mode;     /*  default processing mode  */
object   r;        /*  the window obj           */
@end group
@end example
This method is used to determine when or if the Windows default message
procedure is processed for a given message (@code{msg}) associated with
a particular window (@code{wind}).

WDS allows a programmer to specify an arbitrary number of functions
to be executed whenever a window receives a specific message (via
@code{AddHandlerAfter} and @code{AddHandlerBefore}).  Windows has
default procedures associated with many window messages.  At times
it is necessary to replace or augment this default functionality.
@code{DefaultProcessingMode} gives the programmer control over when and
if this default Windows functionality.  @code{mode} is used to specify
the desired mode.  The following table indicates the valid modes:

@table @code
@item 0
Do not execute the Windows default processing
@item 1
Execute default processing @emph{after} programmer defined handlers
@item 2
Execute default processing @emph{before} programmer defined handlers
@end table

Note that the default mode is always @code{1}, and must be explicitly
changed, if desired, for each message associated with each window.

@code{msg} is the particular message you wish to affect.  These messages
are fully documented in the Windows documentation in the Messages
section.  They normally begin with @code{WM_}.
@example
@group
@exdent Example:

gDefaultProcessingMode(wind, (unsigned) WM_SIZE, 0);
@end group
@end example
@sp 1
See also:  @code{AddHandlerAfter}
@end deffn








@deffn {Disable} Disable::Window
@sp 2
@example
@group
r = gDisable(wind);

object  wind;   /*  a window object  */
object  r;      /*  wind             */
@end group
@end example
This method disables window @code{wind} for user input.  A
disabled window is grayed out and does not accept keyboard or
mouse input.

The window object passed is returned.
@example
@group
@exdent Example:

gDisable(myWind);
@end group
@end example
@sp 1
See also:  @code{Enable}
@end deffn




@deffn {DisableAll} DisableAll::Window
@sp 2
@example
@group
r = gDisableAll(wind);

object  wind;   /*  a window object  */
object  r;      /*  wind             */
@end group
@end example
This method disables all controls contained within window
@code{wind}.

The window object passed is returned.
@example
@group
@exdent Example:

gDisableAll(myWind);
@end group
@end example
@sp 1
See also:  @code{EnableAll, Disable}
@end deffn




@deffn {Dispose} Dispose::Window
@sp 2
@example
@group
r = gDispose(wind);

object  wind;   /*  a window object    */
object  r;      /*  NULL               */
@end group
@end example
This method is used to remove and dispose of a window object when it
is no longer needed.  This method should be called on all windows
when they are no longer needed.

The value returned is always @code{NULL} and may be used to null out
the variable which contained the object being disposed in order to
avoid future accidental use.
@example
@group
@exdent Example:

object  myWind;

myWind = gDispose(myWind);
@end group
@end example
@sp 1
See also:  @code{AutoDispose, New::ChildWindow, New::PopupWindow}
@end deffn










@deffn {DrawText} DrawText::Window
@sp 2
@example
@group
r = gDrawText(wind, str, rect, fmt);

object  wind;   /*  a window object    */
char   *str;    /*  text to draw       */
RECT   *rect;   /*  bounding rectangle */
int     fmt;    /*  format flags       */
object  r;      /*  wind               */
@end group
@end example
This method draws text @code{str} in window @code{wind} at the
position and size specified by @code{rect}.  The @code{fmt}
argument specifies formatting flags as defined by the Windows
@code{DrawText} API (e.g.@: @code{DT_CENTER}, @code{DT_WORDBREAK}).
@example
@group
@exdent Example:

RECT  rc = @{10, 10, 200, 50@};

gDrawText(myWind, "Hello", &rc, DT_CENTER);
@end group
@end example
@c @sp 1
@c See also:  @code{}
@end deffn




@deffn {Enable} Enable::Window
@sp 2
@example
@group
r = gEnable(wind);

object  wind;   /*  a window object  */
object  r;      /*  wind             */
@end group
@end example
This method enables window @code{wind} for user input, reversing
a prior call to @code{Disable}.

The window object passed is returned.
@example
@group
@exdent Example:

gEnable(myWind);
@end group
@end example
@sp 1
See also:  @code{Disable}
@end deffn




@deffn {EnableAll} EnableAll::Window
@sp 2
@example
@group
r = gEnableAll(wind);

object  wind;   /*  a window object  */
object  r;      /*  wind             */
@end group
@end example
This method enables all controls contained within window
@code{wind}.

The window object passed is returned.
@example
@group
@exdent Example:

gEnableAll(myWind);
@end group
@end example
@sp 1
See also:  @code{DisableAll, Enable}
@end deffn




@deffn {Erase} Erase::Window
@sp 2
@example
@group
r = gErase(wind, brow, erow, bcol, ecol);

object  wind;   /*  a window object     */
int     brow;   /*  beginning row       */
int     erow;   /*  ending row          */
int     bcol;   /*  beginning column    */
int     ecol;   /*  ending column       */
object  r;      /*  the window object   */
@end group
@end example
This method is used to delete all text vectors which intersect the
rectangle defined by the parameters.  These parameters are adjusted
according to the current scaling mode (set with
@code{SetScalingMode::Application}.  All coordinates are measured from
the upper left hand corner of the window and start with @code{0,0}.

The object returned is the window object passed.
@example
@group
@exdent Example:

gErase(win, 3, 10, 0, 200);
@end group
@end example
@sp 1
See also:  @code{EraseLines, EraseAll}
@end deffn










@deffn {EraseAll} EraseAll::Window
@sp 2
@example
@group
r = gEraseAll(wind);

object  wind;   /*  a window object     */
object  r;      /*  the window object   */
@end group
@end example
This method is used to erase all text lines associated with a window.

The object returned is the window object passed.
@example
@group
@exdent Example:

gEraseAll(win);
@end group
@end example
@sp 1
See also:  @code{EraseLines, Erase}
@end deffn










@deffn {EraseLines} EraseLines::Window
@sp 2
@example
@group
r = gEraseLines(wind, blin, elin);

object  wind;   /*  a window object     */
int     blin;   /*  beginning line      */
int     elin;   /*  ending line         */
object  r;      /*  the window object   */
@end group
@end example
This method is used to delete all text lines from a starting line number
(@code{blin}) to an ending line number (@code{elin}). All lines are
measured from the upper left hand corner of the window and start with
@code{0,0}.  Lines positions are calculated based on the current font.

The object returned is the window object passed.
@example
@group
@exdent Example:

gEraseLines(win, 3, 10);
@end group
@end example
@sp 1
See also:  @code{Erase, EraseAll}
@end deffn









@deffn {Getch} Getch::Window
@sp 2
@example
@group
ch = gGetch(wind);

object  wind;   /*  a window object   */
int     ch;     /*  character read    */
@end group
@end example
This method is used to obtain the next input character struck by the
user.  The @code{SetBlock} method may be used to determine whether or not
@code{Getch} waits if a character is not available.  The value returned
is the character struck by the user or 0 if no character was available
and the input was non-blocking.
@example
@group
@exdent Example:

object  myWind;
int     ch;

ch = gGetch(myWind);
@end group
@end example
@sp 1
See also:  @code{Kbhit, SetBlock, Gets}
@end deffn
















@deffn {GetParent} GetParent::Window
@sp 2
@example
@group
prnt = gGetParent(wind);

object  wind;   /*  child window object   */
object  prnt;   /*  parent window object  */
@end group
@end example
This method is used to obtain the parent window object associated
with window @code{wind}.
@example
@group
@exdent Example:

object  myWind, parentWind;

parentWind = gGetParent(myWind);
@end group
@end example
@sp 1
See also:  @code{SetParent, New::ChildWindow}
@end deffn















@deffn {GetPosition} GetPosition::Window
@sp 2
@example
@group
r = gGetPosition(wind, vert, horz);

object  wind;   /*  a window object     */
int    *vert;   /*  vertical position   */
int    *horz;   /*  horizontal position */
object  r;      /*  wind arg passed     */
@end group
@end example
This method is used to get the current position of a window.
@code{vert} and @code{horz} are in increments dictated by the
mode selected by @code{SetScalingMode::Application}.  The window
object passed is returned.
@example
@group
@exdent Example:

int     y, x;

gGetPosition(myWind, &y, &x);
@end group
@end example
@sp 1
See also:  @code{SetScalingMode::Application, GetSize, SetPosition}
@end deffn







@deffn {Gets} Gets::Window
@sp 2
@example
@group
r = gGets(wind, buf, len);

object   wind;  /*  a window object       */
char    *buf;   /*  input buffer          */
unsigned len;   /*  length of buffer      */
char    *r;     /*  buf                   */
@end group
@end example
This method is used to read (accept) a single line of text from a user
until a return is hit or @code{len}-1 key strokes have been entered
and place them in @code{buf}.

This method returns a pointer to the buffer passed unless there is an
error, in which case @code{NULL} is returned.
The input accepted will a return terminated line entered by the user
up to @code{len}-1 characters.  However, the number of characters
may be less if non-blocking
io is selected (via @code{SetBlock}) and there aren't enough characters
available.  @code{SetRaw} may also be used to control whether backspace
processing will be performed.
@example
@group
@exdent Example:

object  myWind;
char    buf[80];

gGets(myWind, buf, sizeof(buf));
@end group
@end example
@sp 1
See also:  @code{SetBlock, SetRaw, Gets, Getch}
@end deffn












@deffn {GetSize} GetSize::Window
@sp 2
@example
@group
r = gGetSize(wind, vert, horz);

object  wind;   /*  a window object     */
int    *vert;   /*  vertical size       */
int    *horz;   /*  horizontal size     */
object  r;      /*  wind arg passed     */
@end group
@end example
This method is used to get the current size of a window.
@code{vert} and @code{horz} are in increments dictated by the
mode selected by @code{SetScalingMode::Application}.  The window
object passed is returned.
@example
@group
@exdent Example:

int     y, x;

gGetSize(myWind, &y, &x);
@end group
@end example
@sp 1
See also:  @code{SetScalingMode::Application, GetPosition, SetSize,}
@iftex
@hfil @break @hglue .63in 
@end iftex
@code{GetSize::Application}
@end deffn










@deffn {GetStyle} GetStyle::Window
@sp 2
@example
@group
r = gGetStyle(wind);

object  wind;   /*  a window object   */
DWORD   r;      /*  current style     */
@end group
@end example
This method returns the current Win32 window style associated with
window @code{wind}.  The style values are as defined by the Windows
documentation under @code{CreateWindow}.
@example
@group
@exdent Example:

DWORD  style;

style = gGetStyle(myWind);
@end group
@end example
@sp 1
See also:  @code{SetStyle}
@end deffn




@deffn {GetTag} GetTag::Window
@sp 2
@example
@group
r = gGetTag(wind);

object  wind;   /*  a window object   */
object  r;      /*  tag               */
@end group
@end example
This method is used to obtain a Dynace object which has been associated
with a window via @code{SetTag}.  The value return is the object which has
been associated with the window object @code{wind}.  If there is no object
associated with the window, @code{NULL} will be returned.
@example
@group
@exdent Example:

object  myWind, someObj;

someObj = gGetTag(myWind);
@end group
@end example
@sp 1
See also:  @code{SetTag, SetTag::Dialog}
@end deffn












@deffn {GetTask} GetTask::Window
@sp 2
@example
@group
r = gGetTask(wind);

object  wind;   /*  a window object  */
object  r;      /*  task object      */
@end group
@end example
This method returns the Task object associated with window
@code{wind}, or @code{NULL} if no task has been set.
@example
@group
@exdent Example:

object  task;

task = gGetTask(myWind);
@end group
@end example
@sp 1
See also:  @code{SetTask}
@end deffn




@deffn {GetTopic} GetTopic::Window
@sp 2
@example
@group
r = gGetTopic(wind);

object  wind;   /*  a window object  */
char   *r;      /*  help topic       */
@end group
@end example
This method returns the help topic string associated with window
@code{wind}, or @code{NULL} if no topic has been set.
@example
@group
@exdent Example:

char  *topic;

topic = gGetTopic(myWind);
@end group
@end example
@sp 1
See also:  @code{SetTopic}
@end deffn




@deffn {Handle} Handle::Window
@sp 2
@example
@group
h = gHandle(wind);

object  wind;   /*  window object  */
HANDLE  h;      /*  Windows handle */
@end group
@end example
This method is used to obtain the Windows internal handle associated with
a window object.  Note that this will be @code{0} prior to executing
@code{Show} on the window because that is the point where Windows creates
the handle.

Note that this method may be used with most WDS objects in order to obtain
the internal handle that Windows normally associates with each type of object.
See the appropriate documentation.
@example
@group
@exdent Example:

object  myWind;
HANDLE  h;

h = gHandle(myWind);
@end group
@end example
@c @sp 1
@c See also:  @code{Message}
@end deffn






@deffn {InitFunction} InitFunction::Window
@sp 2
@example
@group
r = gInitFunction(wind, fun);

object  wind;      /*  a window object         */
int    (*fun)();   /*  initialization function  */
object  r;         /*  wind                    */
@end group
@end example
This method sets the initialization function associated with window
@code{wind} to @code{fun}.  The init function is called when the
window is initialized.  @code{fun} receives the window object as
its argument.

Any previously set init function is replaced.

The window object passed is returned.
@example
@group
@exdent Example:

static  int  my_init(object wind)
@{
        /* perform initialization */
        return 0;
@}

        .
        .
        gInitFunction(myWind, my_init);
@end group
@end example
@sp 1
See also:  @code{AddInitFunctionAfter, AddInitFunctionBefore}
@end deffn




@deffn {IsIconic} IsIconic::Window
@sp 2
@example
@group
r = gIsIconic(wind);

object  wind;   /*  a window object  */
int     r;      /*  1 if minimized   */
@end group
@end example
This method returns 1 if window @code{wind} is currently minimized
(iconic), or 0 otherwise.
@example
@group
@exdent Example:

if (gIsIconic(myWind))
        /* window is minimized */
@end group
@end example
@sp 1
See also:  @code{IsMaximized, IsNormal}
@end deffn




@deffn {IsMaximized} IsMaximized::Window
@sp 2
@example
@group
r = gIsMaximized(wind);

object  wind;   /*  a window object  */
int     r;      /*  1 if maximized   */
@end group
@end example
This method returns 1 if window @code{wind} is currently maximized,
or 0 otherwise.
@example
@group
@exdent Example:

if (gIsMaximized(myWind))
        /* window is maximized */
@end group
@end example
@sp 1
See also:  @code{IsIconic, IsNormal}
@end deffn




@deffn {IsNormal} IsNormal::Window
@sp 2
@example
@group
r = gIsNormal(wind);

object  wind;   /*  a window object    */
int     r;      /*  1 if normal state  */
@end group
@end example
This method returns 1 if window @code{wind} is in the normal
(restored) state --- that is, neither minimized nor maximized.
@example
@group
@exdent Example:

if (gIsNormal(myWind))
        /* window is in normal state */
@end group
@end example
@sp 1
See also:  @code{IsIconic, IsMaximized}
@end deffn




@deffn {Kbhit} Kbhit::Window
@sp 2
@example
@group
r = gKbhit(wind);

object  wind;   /*  a window object   */
int     r;      /*  characters ready  */
@end group
@end example
This method is used to determine if and how many characters are waiting
in the character input queue.  This method is used in an attempt to
provide a familiar character based user input mechanism.

The value returned is the number of characters ready for input or
zero if none are available.
@example
@group
@exdent Example:

object  myWind;
int     n;

n = gKbhit(myWind);
@end group
@end example
@sp 1
See also:  @code{Getch, SetBlock, Gets}
@end deffn



















@deffn {KillTimer} KillTimer::Window
@sp 2
@example
@group
r = gKillTimer(wind, timerID);

object  wind;      /*  a window object  */
int     timerID;   /*  timer ID         */
int     r;         /*  result           */
@end group
@end example
This method removes a timer previously created with @code{AddTimer}.
@code{timerID} is the value returned by @code{AddTimer}.
@example
@group
@exdent Example:

gKillTimer(myWind, timerID);
@end group
@end example
@sp 1
See also:  @code{AddTimer}
@end deffn




@deffn {LoadCursor} LoadCursor::Window
@sp 2
@example
@group
r = mLoadCursor(wind, csr);

object   wind;  /*  a window object    */
unsigned csr;   /*  cursor identifier  */
object   r;     /*  cursor object      */
@end group
@end example
This method is used to load a programmer defined cursor and associate it
with window @code{wind}.  The cursor would then be displayed any time
the pointer was placed in the window.

@code{csr} is a programmer defined unsigned integer which identifies
the cursor.  This identifier is normally a macro and defined through the
resource editor.  The cursor object will be automatically destroyed
whenever the window is destroyed or if a new cursor is loaded.

The value returned is an object representing the cursor loaded, or
@code{NULL} if the cursor was not found.
@example
@group
@exdent Example:

object  myWind;

mLoadCursor(myWind, MY_CURSOR);
@end group
@end example
@sp 1
See also:  @code{LoadSystemCursor, Load::ExternalCursor, Use}
@end deffn











@deffn {LoadFont} LoadFont::Window
@sp 2
@example
@group
r = gLoadFont(wind, fname, sz);

object   wind;  /*  a window object  */
char    *fname; /*  font name        */
int      sz;    /*  point size       */
object   r;     /*  font object      */
@end group
@end example
This method is used to load an arbitrary font by name at any point size
and associate it as the default font for future text output to window
@code{wind}.  @code{fname} is the full name of the font as it appears
when you list the available fonts via the control-panel / fonts Windows
utility, minus the font type in parentheses.  @code{sz} indicates the
desired point size.

All font objects associated with a window will be automatically
destroyed whenever the window is destroyed.

The value returned is an object representing the font loaded, or
@code{NULL} if the font was not found.

Note that @code{fname} may also be a Dynace object cast as a (@code{char *}).
@example
@group
@exdent Example:

object  myWind;

gLoadFont(myWind, "Times New Roman", 12);
@end group
@end example
@sp 1
See also:  @code{LoadSystemFont, Indirect::ExternalFont, Use}
@end deffn









@deffn {LoadIcon} LoadIcon::Window
@sp 2
@example
@group
r = mLoadIcon(wind, icn);

object   wind;  /*  a window object  */
unsigned icn;   /*  icon identifier  */
object   r;     /*  icon object      */
@end group
@end example
This method is used to load a programmer defined icon and associate it
with window @code{wind}.  The icon would then be displayed if window
@code{wind} was iconized.

@code{icn} is a programmer defined unsigned integer which identifies
the icon.  This identifier is normally a macro and defined through the
resource editor.  The icon object will be automatically destroyed
whenever the window is destroyed or if a new icon is loaded.

The value returned is an object representing the icon loaded, or
@code{NULL} if the icon was not found.
@example
@group
@exdent Example:

object  myWind;

mLoadIcon(myWind, ALGOCORP_ICON);
@end group
@end example
@sp 1
See also:  @code{LoadSystemIcon, Load::ExternalIcon, Use}
@end deffn








@deffn {LoadMenu} LoadMenu::Window
@sp 2
@example
@group
r = mLoadMenu(wind, mnu);

object   wind;  /*  a window object  */
unsigned mnu;   /*  menu identifier  */
object   r;     /*  menu object      */
@end group
@end example
This method is used to load a programmer defined menu and associate it
with window @code{wind}.  The menu would be displayed at the top
of the window.

@code{mnu} is a programmer defined unsigned integer which identifies the
menu.  This identifier is normally a macro and defined through the
resource editor.  Any menu previously associated with the window is
pushed on a stack so that previous menus may be easily returned to in a
last-in-first-out basis using the @code{PopMenu} method.  All menus
associated with a window will be destroyed when the window is destroyed.

The value returned is an object representing the menu loaded, or
@code{NULL} if the menu was not found.
@example
@group
@exdent Example:

object  myWind;

mLoadMenu(myWind, MY_MENU);
@end group
@end example
@sp 1
See also:  @code{Associate, PopMenu, LoadMenuStr, Load::ExternalMenu, Use}
@end deffn














@deffn {LoadMenuStr} LoadMenuStr::Window
@sp 2
@example
@group
r = gLoadMenuStr(wind, mnu);

object   wind;  /*  a window object  */
char    *mnu;   /*  menu identifier  */
object   r;     /*  menu object      */
@end group
@end example
This method is used to load a programmer defined menu and associate it
with window @code{wind}.  The menu would be displayed at the top
of the window.

@code{mnu} is a programmer defined name which identifies the menu.  This
name is defined through the resource editor.  Any menu previously associated
with the window is pushed on a stack so that previous menus may be easily
returned to in a last-in-first-out basis using the @code{PopMenu} method.
All menus associated with a window will be destroyed when the window
is destroyed.

Note that the @code{LoadMenu} method is used more frequently since the
resource editors assign macros to each menu name.

The value returned is an object representing the menu loaded, or
@code{NULL} if the menu was not found.
@example
@group
@exdent Example:

object  myWind;

gLoadMenuStr(myWind, "mymenu");
@end group
@end example
@sp 1
See also:  @code{Associate, LoadMenu, LoadStr::ExternalMenu, Use}
@end deffn












@deffn {LoadSystemCursor} LoadSystemCursor::Window
@sp 2
@example
@group
r = gLoadSystemCursor(wind, csr);

object   wind;  /*  a window object    */
LPCSTR   csr;   /*  cursor identifier  */
object   r;     /*  cursor object      */
@end group
@end example
This method is used to load a Windows predefined cursor and associate it
with window @code{wind}.  The cursor would then be displayed if the pointer
was positioned over the window.

@code{csr} is a Windows defined macro which identifies the cursor.  The
available options are defined in the Windows documentation under the
function named @code{LoadCursor} and normally begin with @code{IDC_}.  The
cursor object will be automatically destroyed whenever the window is
destroyed or if a new cursor is loaded.

The value returned is an object representing the cursor loaded, or
@code{NULL} if the cursor was not found.
@example
@group
@exdent Example:

object  myWind;

gLoadSystemCursor(myWind, IDC_CROSS);
@end group
@end example
@sp 1
See also:  @code{LoadCursor, LoadSys::SystemCursor, Use}
@end deffn











@deffn {LoadSystemFont} LoadSystemFont::Window
@sp 2
@example
@group
r = gLoadSystemFont(wind, fnt);

object   wind;  /*  a window object  */
unsigned fnt;   /*  font identifier  */
object   r;     /*  font object      */
@end group
@end example
This method is used to load a Windows predefined font and associate it
with window @code{wind}.  The last font associated with a window is the
one which will be used when any text is output to the window.

@code{fnt} is a Windows defined macro which identifies the font.  The
available options are defined in the Windows documentation under the
function named @code{GetStockObject} and normally end with @code{_FONT}.
The font object will be automatically destroyed whenever the window is
destroyed.

The value returned is an object representing the font, or
@code{NULL} if the font was not found.
@example
@group
@exdent Example:

object  myWind;

gLoadSystemFont(myWind, SYSTEM_FONT);
@end group
@end example
@sp 1
See also:  @code{LoadFont, Load::SystemFont, Use}
@end deffn
















@deffn {LoadSystemIcon} LoadSystemIcon::Window
@sp 2
@example
@group
r = gLoadSystemIcon(wind, icn);

object   wind;  /*  a window object  */
LPCSTR   icn;   /*  icon identifier  */
object   r;     /*  icon object      */
@end group
@end example
This method is used to load a Windows predefined icon and associate it
with window @code{wind}.  The icon would then be displayed if window
@code{wind} was iconized.

@code{icn} is a Windows defined macro which identifies the icon.  The
available options are defined in the Windows documentation under the
function named @code{LoadIcon} and normally begin with @code{IDI_}.  The
icon object will be automatically destroyed whenever the window is
destroyed or if a new icon is loaded.

The value returned is an object representing the icon loaded, or
@code{NULL} if the icon was not found.
@example
@group
@exdent Example:

object  myWind;

gLoadSystemIcon(myWind, IDI_APPLICATION);
@end group
@end example
@sp 1
See also:  @code{LoadIcon, LoadSys::SystemIcon, Use}
@end deffn














@deffn {Maximize} Maximize::Window
@sp 2
@example
@group
r = gMaximize(wind);

object  wind;   /*  a window object  */
object  r;      /*  wind             */
@end group
@end example
This method maximizes window @code{wind} so that it fills the
entire screen.

The window object passed is returned.
@example
@group
@exdent Example:

gMaximize(myWind);
@end group
@end example
@c @sp 1
@c See also:  @code{}
@end deffn




@deffn {Menu} Menu::Window
@sp 2
@example
@group
r = gMenu(wind);

object  wind;   /*  a window object  */
object  r;      /*  menu object      */
@end group
@end example
This method returns the menu object associated with window
@code{wind}, or @code{NULL} if no menu has been loaded.
@example
@group
@exdent Example:

object  menu;

menu = gMenu(myWind);
@end group
@end example
@sp 1
See also:  @code{LoadMenu, LoadMenuStr}
@end deffn




@deffn {MenuItemMode} MenuItemMode::Window
@sp 2
@example
@group
r = mMenuItemMode(wind, itm, mod);

object  wind;   /*  a window object     */
unsigned itm;   /*  menu item           */
unsigned mod;   /*  menu item mode      */
object   r;        /*  the menu         */
@end group
@end example
This method is used to set the mode associated with a particular item
in the menu which is currently attached to window @code{wind}.

@code{itm} is a programmer defined macro which identifies one particular
choice among those available within the menu which is currently attached
to window @code{wind}.  This macro is defined while the programmer defines
the entire menu using the resource editor.

@code{mod} may be one of @code{MF_DISABLED}, @code{MF_ENABLED} or
@code{MF_GRAYED} and is documented in the Windows documentation under
the function @code{EnableMenuItem}.

The menu object associated with the window is returned.
@example
@group
@exdent Example:

mMenuItemMode(win, ID_FILE_MESSAGE, MF_GRAYED);
@end group
@end example
@sp 1
See also:  @code{LoadMenu, Associate}
@end deffn











@deffn {Message} Message::Window
@sp 2
@example
@group
r = gMessage(wind, msg);

object  wind;   /*  window object  */
char    *msg;   /*  message        */
object  r;      /*  window object  */
@end group
@end example
This method is used to open up a temporary informational window.  The
window will contain the message given by @code{msg} and the user must
acknowledge the window prior to continuing by hitting an OK button.


The value returned is the window passed.
@example
@group
@exdent Example:

object  myWind;

gMessage(myWind, "Press OK to continue.");
@end group
@end example
@sp 1
See also:  @code{MessageWithTopic}
@end deffn












@deffn {MessageWithTopic} MessageWithTopic::Window
@sp 2
@example
@group
r = gMessageWithTopic(wind, msg, tpc);

object  wind;   /*  window object  */
char    *msg;   /*  message        */
char    *tpc;   /*  help topic     */
object  r;      /*  window object  */
@end group
@end example
This method is used to open up a temporary informational window.  The
window will contain the message given by @code{msg} and the user must
acknowledge the window prior to continuing by hitting an OK button.

If the user hits the @code{F1} key while presented with the message,
the help topic identified by @code{tpc} will get displayed via the Windows
help system.

The value returned is the window passed.
@example
@group
@exdent Example:

object  myWind;

gMessageWithTopic(myWind, "Press OK to continue.", "mytopic");
@end group
@end example
@sp 1
See also:  @code{Message} and the @code{HelpSystem} class.
@end deffn











@deffn {MoveToTop} MoveToTop::Window
@sp 2
@example
@group
r = gMoveToTop(wind);

object  wind;   /*  a window object  */
object  r;      /*  wind             */
@end group
@end example
This method brings window @code{wind} to the top of the Z-order,
making it the foreground window.

The window object passed is returned.
@example
@group
@exdent Example:

gMoveToTop(myWind);
@end group
@end example
@sp 1
See also:  @code{SetZOrder}
@end deffn




@deffn {New} New::Window
@sp 2
@example
@group
r = vNew(Window);

object  r;      /*  new window  */
@end group
@end example
This class method is used to create a basic window object.  It is used
by all the subclasses of @code{Window} and would not normally be
used by a programmer.

The value returned is the new window created.
@example
@group
@exdent Example:

object  myWind;

myWind = vNew(Window);
@end group
@end example
@sp 1
See also:  @code{New::MainWindow, New::ChildWindow, New::PopupWindow, Show}
@iftex
@hfil @break @hglue .63in   
@end iftex
@code{Dispose}
@end deffn







@deffn {NewBuiltIn} NewBuiltIn::Window
@sp 2
@example
@group
r = gNewBuiltIn(Window, class, parent);

char    *class; /*  Windows built in class designation  */
object  parent; /*  parent window object                */
object  r;      /*  new child window                    */
@end group
@end example
This class method is used to create a child window which is used as a
Windows-built-in control.  @code{Window} represents the @code{Window}
class and is typed in as shown.  @code{class} is a string representing
one of the window classes build into Windows.  This string is defined
and documented by Windows under the function named @code{CreateWindow}.
@code{parent} represents the parent window object and must be specified.

The value returned is the new child window (control) created.
@example
@group
@exdent Example:

object  myWind, ctl;

myWind = vNew(MainWindow, "App Name");
ctl = gNewBuiltIn(Window, "button", myWind);
@end group
@end example
@sp 1
See also:  @code{New::ButtonWindow}
@end deffn












@deffn {NextControl} NextControl::Window
@sp 2
@example
@group
r = gNextControl(wind, ctl);

object  wind;   /*  a window object      */
object  ctl;    /*  current control      */
object  r;      /*  next control or NULL */
@end group
@end example
This method returns the next focusable control in window @code{wind}
after control @code{ctl}.  If @code{ctl} is @code{NULL}, the first
focusable control in the window is returned.
@example
@group
@exdent Example:

object  next;

next = gNextControl(myWind, NULL);  /* get first */
@end group
@end example
@sp 1
See also:  @code{PreviousControl}
@end deffn




@deffn {Perform} Perform::Window
@sp 2
@example
@group
r = gPerform(wind);

object  wind;   /*  a window object  */
int     r;      /*  result           */
@end group
@end example
This method sets up controls and shows window @code{wind} in a
non-blocking manner.  Control returns to the caller immediately.
The window remains visible and active until closed by the user or
programmatically.
@example
@group
@exdent Example:

gPerform(myWind);
@end group
@end example
@sp 1
See also:  @code{PerformModal, Show}
@end deffn




@deffn {PerformModal} PerformModal::Window
@sp 2
@example
@group
r = gPerformModal(wind);

object  wind;   /*  a window object  */
int     r;      /*  result           */
@end group
@end example
This method shows window @code{wind} modally, blocking the caller
until the window is closed.
@example
@group
@exdent Example:

int  result;

result = gPerformModal(myWind);
@end group
@end example
@sp 1
See also:  @code{Perform}
@end deffn




@deffn {PopMenu} PopMenu::Window
@sp 2
@example
@group
r = gPopMenu(wind);

object   wind;  /*  a window object  */
object   r;     /*  menu object      */
@end group
@end example
This method is used to remove and destroy the current menu associated
with window @code{wind} and restore the previous menu associated with
the window.  A last-in-first-out list of menus associated with a window
may be established via the @code{LoadMenu}, @code{LoadMenuStr} or
@code{Use} methods.

All function associations and modes previously associated with the
new menu will be restored.

The object returned is the new menu object.
@example
@group
@exdent Example:

object  myWind;

gPopMenu(myWind);
@end group
@end example
@sp 1
See also:  @code{LoadMenu, Use}
@end deffn











@deffn {PreviousControl} PreviousControl::Window
@sp 2
@example
@group
r = gPreviousControl(wind, ctl);

object  wind;   /*  a window object          */
object  ctl;    /*  current control          */
object  r;      /*  previous control or NULL */
@end group
@end example
This method returns the previous focusable control in window
@code{wind} before control @code{ctl}.  If @code{ctl} is
@code{NULL}, the last focusable control in the window is returned.
@example
@group
@exdent Example:

object  prev;

prev = gPreviousControl(myWind, currentCtl);
@end group
@end example
@sp 1
See also:  @code{NextControl}
@end deffn




@deffn {Printf} Printf::Window
@sp 2
@example
@group
r = vPrintf(wind, fmt, ...);

object   wind;  /*  a window object       */
char    *fmt;   /*  format string         */
int      r;     /*  length of output      */
@end group
@end example
This method is used to display a string of text (@code{str}) in a sequential
fashion on window @code{wind}.  It is analogous to the standard C library
function @code{fprintf}.  All of the standard features of your C library
@code{fprintf} function are supported and that documentation should be
consulted for full documentation on the arguments.

This method is used to support the standard streams interface, is
actually defined by the @code{Stream} class, and is documented here for
convenience.

Note that this method begins with ``v'' since it takes variable arguments.
The length of the resulting output will be returned.
@example
@group
@exdent Example:

object  myWind;
int     age = 32;

vPrintf(myWind, "My age =- %d\n", age);
@end group
@end example
@sp 1
See also:  @code{TextOut, Puts, Write}
@end deffn



















@deffn {PrintScreen} PrintScreen::Window
@sp 2
@example
@group
r = gPrintScreen(wind, printerObj);

object  wind;        /*  a window object           */
object  printerObj;  /*  printer object or NULL     */
object  r;           /*  wind                      */
@end group
@end example
This method prints the contents of window @code{wind} to the printer
specified by @code{printerObj}.  If @code{printerObj} is @code{NULL},
the default printer is used.

The window object passed is returned.
@example
@group
@exdent Example:

gPrintScreen(myWind, NULL);  /* print to default printer */
@end group
@end example
@c @sp 1
@c See also:  @code{}
@end deffn




@deffn {Puts} Puts::Window
@sp 2
@example
@group
r = gPuts(wind, str);

object   wind;  /*  a window object       */
char    *str;   /*  string to be output   */
int      r;     /*  length of str         */
@end group
@end example
This method is used to display a string of text (@code{str}) in a sequential
fashion on window @code{wind}.  This method is used to support the standard
streams interface, is actually defined by the @code{Stream} class, and
is documented here for convenience.

@code{str} is the text to be displayed. Text @code{str} will be
displayed on window @code{wind} and its length will be returned.
@example
@group
@exdent Example:

object  myWind;

gPuts(myWind, "Hello World\n");
@end group
@end example
@sp 1
See also:  @code{TextOut, Printf, Write}
@end deffn











@deffn {Read} Read::Window
@sp 2
@example
@group
r = gRead(wind, buf, len);

object   wind;  /*  a window object       */
char    *buf;   /*  input buffer          */
unsigned len;   /*  length of buffer      */
int      r;     /*  bytes read            */
@end group
@end example
This method is used to read (accept) @code{len} key strokes from the user
and place them in @code{buf}.

Since the @code{Window} class is a subclass of @code{Stream}, this
method is mainly provided to support the standard interface dictated by
the @code{Stream} class.

This method returns the number of characters actually accepted.  It will
not be more than @code{len}, but may be less if non-blocking io is selected
(via @code{SetBlock}) and there aren't enough characters available.
@code{SetRaw} may also be used to control whether backspace processing will
be performed.
@example
@group
@exdent Example:

object  myWind;
char    buf[80];

gRead(myWind, buf, sizeof(buf)-1);
@end group
@end example
@sp 1
See also:  @code{SetBlock, SetRaw, Gets, Getch}
@end deffn














@deffn {RedrawMenu} RedrawMenu::Window
@sp 2
@example
@group
r = gRedrawMenu(wind);

object  wind;   /*  a window object  */
object  r;      /*  wind             */
@end group
@end example
This method redraws the menu bar of window @code{wind}.  This is
useful after programmatically modifying menu items.

The window object passed is returned.
@example
@group
@exdent Example:

gRedrawMenu(myWind);
@end group
@end example
@sp 1
See also:  @code{RedrawWindow, LoadMenu}
@end deffn




@deffn {RedrawWindow} RedrawWindow::Window
@sp 2
@example
@group
r = gRedrawWindow(wind);

object  wind;   /*  a window object  */
object  r;      /*  wind             */
@end group
@end example
This method invalidates the entire client area of window @code{wind},
causing Windows to send a repaint message.

The window object passed is returned.
@example
@group
@exdent Example:

gRedrawWindow(myWind);
@end group
@end example
@sp 1
See also:  @code{RedrawMenu}
@end deffn




@deffn {ScrollHorz} ScrollHorz::Window
@sp 2
@example
@group
r = gScrollHorz(wind, cols);

object  wind;   /*  a window object     */
int     cols;   /*  columns to scroll   */
object  r;      /*  the window object   */
@end group
@end example
This method is used to perform a horizontal scrolling of the text in
window @code{wind}.  This has the same effect as if the user caused
horizontal scrolling via the horizontal scroll bar at the bottom
of the window.  No text is lost, a different portion of the text is
displayed.

@code{wind} is the window which is to be affected, and @code{cols} is the
number of columns to scroll.  If @code{cols} is positive the scroll moves
the text to the left, and negative moves the text to the right.

The scaling factor associated with @code{cols} is set by
@code{SetScalingMode::Application}.
@example
@group
@exdent Example:

object  myWind;

gScrollHorz(myWind, 2);
@end group
@end example
@sp 1
See also:  @code{ScrollVert, SetScalingMode::Application}
@end deffn











@deffn {ScrollVert} ScrollVert::Window
@sp 2
@example
@group
r = gScrollVert(wind, rows);

object  wind;   /*  a window object     */
int     rows;   /*  rows to scroll      */
object  r;      /*  the window object   */
@end group
@end example
This method is used to perform a vertical scrolling of the text in
window @code{wind}.  This has the same effect as if the user caused
vertical scrolling via the vertical scroll bar at the right side
of the window.  No text is lost, a different portion of the text is
displayed.

@code{wind} is the window which is to be affected, and @code{rows} is the
number of rows to scroll.  If @code{rows} is positive the scroll moves
the text up, and negative moves the text down.

The scaling factor associated with @code{rows} is set by
@code{SetScalingMode::Application}.
@example
@group
@exdent Example:

object  myWind;

gScrollVert(myWind, 2);
@end group
@end example
@sp 1
See also:  @code{ScrollHorz, VertShift, SetScalingMode::Application}
@end deffn










@deffn {SetBlock} SetBlock::Window
@sp 2
@example
@group
r = gSetBlock(wind, flag);

object  wind;   /*  a window object     */
int     flag;   /*  enable/disable flag */
int     r;      /*  previous flag       */
@end group
@end example
This method is used to set the blocking mode associated with keyboard
IO.  It only has effect on @code{Read}, @code{Gets}, and @code{Getch}.
If flag is 1, blocking is enabled, and 0 disables blocking.  When a
new window is created it defaults to blocking enabled.

If blocking is turned on and a keyboard entry is requested (via
@code{Read}, @code{Gets} or @code{Getch}), the keyboard entry function
will not return until the request can be satisfied.  If, however,
blocking is disabled, a keyboard request is made, and insufficient
characters are available, then the input function will return immediately
with a return value indicating the result was short.

The value returned by this method is the previous blocking mode.
@example
@group
@exdent Example:

object  myWind;

gSetBlock(myWind, 0);
@end group
@end example
@sp 1
See also:  @code{Getch, SetRaw, Gets}
@end deffn















@deffn {SetCursor} SetCursor::Window
@sp 2
@example
@group
r = gSetCursor(wind, cursor);

object  wind;     /*  a window object      */
object  cursor;   /*  cursor object or NULL */
object  r;        /*  wind                 */
@end group
@end example
This method sets a custom cursor for window @code{wind}.
Pass @code{NULL} for @code{cursor} to restore the default cursor.

The window object passed is returned.
@example
@group
@exdent Example:

gSetCursor(myWind, myCursor);
gSetCursor(myWind, NULL);  /* restore default */
@end group
@end example
@sp 1
See also:  @code{WaitCursor, LoadCursor}
@end deffn




@deffn {SetFocus} SetFocus::Window
@sp 2
@example
@group
r = gSetFocus(wind);

object  wind;   /*  a window object  */
object  r;      /*  wind             */
@end group
@end example
This method sets the keyboard input focus to window @code{wind}.

The window object passed is returned.
@example
@group
@exdent Example:

gSetFocus(myWind);
@end group
@end example
@c @sp 1
@c See also:  @code{}
@end deffn




@deffn {SetMaxLines} SetMaxLines::Window
@sp 2
@example
@group
r = gSetMaxLines(wind, rows);

object  wind;   /*  a window object     */
int     rows;   /*  max rows            */
int     r;      /*  previous max rows   */
@end group
@end example
This method is used to set or obtain the maximum number of lines of text
which may be associated to a window.  This includes all lines associated
with a window, including lines not being displayed because they are
scrolled off the screen.  Whenever @code{rows} number of lines are
exceeded, WDS automatically and permanently removes the lines at the
logical top of the internal buffer in order to make room for new
lines.

WDS keeps a buffer which holds a programmer definable number of lines of
text.  The user is then able to scroll through this text.  If the
application attempts to display more lines than this maximum (via
@code{vPrintf} for example) the system will automatically call
@code{gVertShift} in order to eliminate the top line and make room for
the new line being appended.

@code{wind} is the window which is to be affected, and @code{rows} is the
maximum number of rows to retain.  If @code{rows} is zero or negative,
the value associated with the window will not be changed.  This is used
to obtain the current value without changing it.

The value returned is the previous value associated with the window.
@example
@group
@exdent Example:

object  myWind;

gSetMaxLines(myWind, 60);
@end group
@end example
@sp 1
See also:  @code{VertShift}
@end deffn













@deffn {SetParent} SetParent::Window
@sp 2
@example
@group
r = gSetParent(wind, prnt);

object  wind;   /*  child window object   */
object  prnt;   /*  parent window object  */
object  r;      /*  wind                  */
@end group
@end example
This method is used to create a child / parent window relationship.
This relationship is automatically established when a child window is
created.  However, this method is provided for increased flexibility.

Whenever the parent window is disposed, WDS will automatically dispose
of all its children windows.  This relationship should normally be
established prior to the child window being @code{Shown}.

The value returned is the child window passed.
@example
@group
@exdent Example:

object  myWind, parentWind;

gSetParent(myWind, parentWind);
@end group
@end example
@sp 1
See also:  @code{GetParent}
@end deffn
















@deffn {SetPosition} SetPosition::Window
@sp 2
@example
@group
r = gSetPosition(wind, vert, horz);

object  wind;   /*  a window object     */
int     vert;   /*  vertical position   */
int     horz;   /*  horizontal position */
object  r;      /*  wind arg passed     */
@end group
@end example
This method is used to set the initial position of a window.
@code{vert} and @code{horz} are in increments dictated by the
mode selected by @code{SetScalingMode::Application}.  The window
object passed is returned.

If this function is not called, Windows will automatically set it to a
reasonable default.
@example
@group
@exdent Example:

object  myWind;

myWind = vNew(MainWindow, "App Name");
gSetPosition(myWind, 3, 10);
@end group
@end example
@sp 1
See also:  @code{SetScalingMode::Application, SetSize, GetPosition}
@end deffn











@deffn {SetRaw} SetRaw::Window
@sp 2
@example
@group
r = gSetRaw(wind, flag);

object  wind;   /*  a window object     */
int     flag;   /*  enable/disable flag */
int     r;      /*  previous flag       */
@end group
@end example
This method is used to set the raw mode associated with keyboard
IO.  It only has effect on @code{Read}, @code{Gets}, and @code{Getch}.
If flag is 1, raw mode is enabled, and 0 disables raw mode.  The
default is raw mode disabled.

Normally, with raw mode disabled, when a user enters keyboard data,
they are able to correct mistakes by using the backspace and reentering
the correct data.  The resulting data the program receives is the
final, corrected input.

When raw mode is enabled, every key the user hits gets returned.  This
includes backspaces.  Therefore, with raw mode enabled, if the user types
``ABD'' followed by backspace and then ``C'', the program will receive all
five characters hit.  However, with raw mode disabled, the program would
only receive the resulting string, ``ABC''.

The value returned by this method is the previous raw mode.
@example
@group
@exdent Example:

object  myWind;

gSetRaw(myWind, 1);
@end group
@end example
@sp 1
See also:  @code{Getch, SetBlock, Gets}
@end deffn











@deffn {SetSize} SetSize::Window
@sp 2
@example
@group
r = gSetSize(wind, vert, horz);

object  wind;   /*  a window object     */
int     vert;   /*  vertical size       */
int     horz;   /*  horizontal size     */
object  r;      /*  wind arg passed     */
@end group
@end example
This method is used to set the initial size of a window.
@code{vert} and @code{horz} are in increments dictated by the
mode selected by @code{SetScalingMode::Application}.  The window
object passed is returned.

If this function is not called, Windows will automatically set it to a
reasonable default.
@example
@group
@exdent Example:

object  myWind;

myWind = vNew(MainWindow, "App Name");
gSetSize(myWind, 10, 40);
@end group
@end example
@sp 1
See also:  @code{SetScalingMode::Application, SetPosition, GetSize}
@end deffn









@deffn {SetStyle} SetStyle::Window
@sp 2
@example
@group
r = gSetStyle(wind, sty);

object  wind;   /*  a window object             */
DWORD   sty;    /*  the window style to use     */
object  r;      /*  returns wind                */
@end group
@end example
This method is used to set the window style associated with window
@code{wind} to style @code{sty}.  The @code{DWORD} data type and valid 
styles are defined by Windows and fully documented by the Windows
documentation.  See the Windows documentation for the function called
@code{CreateWindow}.  The style types normally begin with @code{WS_}
and would be or'ed together to form the selected style.
Note that WDS automatically assigns reasonable defaults to a window
style when it is created.
@example
@group
@exdent Example:

object  myWind;

myWind = vNew(MainWindow, "App Name");
gSetStyle(myWind, WS_OVERLAPPEDWINDOW | WS_VSCROLL | WS_HSCROLL);
@end group
@end example
@c @sp 1
@c See also:  @code{}
@end deffn








@deffn {SetTag} SetTag::Window
@sp 2
@example
@group
r = gSetTag(wind, tag);

object  wind;   /*  a window object   */
object  tag;    /*  tag               */
object  r;      /*  previous tag      */
@end group
@end example
This method is used to associate an arbitrary Dynace object with a
window object.  This may later be retrieved via the @code{GetTag} method.
Since WDS passes around the window object to all @code{Window} methods
this mechanism may be used to pass additional information with the
window.  And since Dynace treats all objects in a uniform manner, this
information attached to the window may be arbitrarily complex.

WDS does not dispose of the tag when the window object is disposed.
This method returns any previous object associated with the window or
@code{NULL}.
@example
@group
@exdent Example:

object  myWind;

gSetTag(myWind, gNewWithInt(ShortInteger, 17));
@end group
@end example
@sp 1
See also:  @code{GetTag, SetTag::Dialog}
@end deffn












@deffn {SetTask} SetTask::Window
@sp 2
@example
@group
r = gSetTask(wind, task);

object  wind;   /*  a window object  */
object  task;   /*  task object      */
object  r;      /*  wind             */
@end group
@end example
This method associates a Task object @code{task} with window
@code{wind}.

The window object passed is returned.
@example
@group
@exdent Example:

gSetTask(myWind, myTask);
@end group
@end example
@sp 1
See also:  @code{GetTask}
@end deffn




@deffn {SetTitle} SetTitle::Window
@sp 2
@example
@group
r = gSetTitle(wind, title);

object  wind;    /*  a window object     */
char    *title;  /*  the new title       */
object  r;       /*  wind arg passed     */
@end group
@end example
This method is used to set the title displayed at the top of window
@code{wind}, if it has a title bar.  @code{title} may also be a Dynace
string @code{object} typecast to a (@code{char *}).  The value returned
is the window (@code{wind}) object passed.
@example
@group
@exdent Example:

object  myWind;

myWind = vNew(MainWindow, "Old Title");
gSetTitle(myWind, "New Title");
@end group
@end example
@sp 1
See also:  @code{New::MainWindow}
@end deffn













@deffn {SetTopic} SetTopic::Window
@sp 2
@example
@group
pt = gSetTopic(wind, tpc);

object  wind;   /*  child window object   */
char    *tpc;   /*  help topic            */
char    *pt;    /*  previous help topic   */
@end group
@end example
This method is used to associate help text with window @code{wind}.
The help text is defined using the Windows help system and labeled
with the topic indicated by @code{tpc}.  Then, if the user hits
the F1 key while in the window, WDS will automatically bring up
the Windows help system and find the indicated topic.

WDS also supports dialog and control specific topics.  See the appropriate
sections.

This method returns any previous topic associated with the window.
@example
@group
@exdent Example:

object  myWind;

gSetTopic(myWind, "myWindHelp");
@end group
@end example
@sp 1
See also:  The @code{HelpSystem} class.
@end deffn













@deffn {SetZOrder} SetZOrder::Window
@sp 2
@example
@group
r = gSetZOrder(wind, m);

object  wind;   /*  a window object     */
HWND    m;      /*  Z-order position    */
object  r;      /*  wind                */
@end group
@end example
This method sets the Z-order position of window @code{wind}.
@code{m} specifies the position and can be @code{HWND_TOP},
@code{HWND_TOPMOST}, @code{HWND_BOTTOM}, or @code{HWND_NOTOPMOST}
as defined by the Windows API.

The window object passed is returned.
@example
@group
@exdent Example:

gSetZOrder(myWind, HWND_TOPMOST);
@end group
@end example
@sp 1
See also:  @code{MoveToTop}
@end deffn




@deffn {Show} Show::Window
@sp 2
@example
@group
r = gShow(wind);

object  wind;   /*  a window object     */
int     r;      /*  always 0            */
@end group
@end example
Once a window object is created (via @code{vNew}) and its various attributes
set, @code{gShow} is called in order to actually create the Windows window,
with the selected attributes, and display (show) it.
@example
@group
@exdent Example:

object  myWind;

myWind = vNew(PopupWindow, "Name", 10, 40);
gShow(myWind);
@end group
@end example
@sp 1
See also:  @code{AutoShow, ProcessMessages::MainWindow}
@end deffn









@deffn {TextBrush} TextBrush::Window
@sp 2
@example
@group
r = gTextBrush(wind, brsh);

object  wind;   /*  a window object   */
object  brsh;   /*  brush object      */
object  r;      /*  wind              */
@end group
@end example
This method is used to determine what brush object is used for foreground
text which is displayed.  Any previously associated brush object will
be disposed.  This brush object will also be automatically disposed when the
window is disposed.

The window passed is returned.
@example
@group
@exdent Example:

object  myWind;

gTextBrush(myWind, vNew(SolidBrush, 255, 0, 0));
@end group
@end example
@sp 1
See also:  @code{BackBrush, Use, SetBackBrush::Application}
        and the @code{Brush} classes
@end deffn










@deffn {TextOut} TextOut::Window
@sp 2
@example
@group
r = gTextOut(wind, row, col, txt);

object  wind;   /*  a window object          */
int     row;    /*  row number of output     */
int     col;    /*  column number of output  */
char    *txt;   /*  string to be output      */
object  r;      /*  wind                     */
@end group
@end example
This method is used to display a string of text (@code{txt}) at row
@code{row} and column @code{col}.  @code{row} and @code{col} have their
origin in the upper left hand corner of the window and are scaled as
dictated by the @code{SetScalingMode::Application} method.
Their index origin is 0.

This method returns the window argument passed.
@example
@group
@exdent Example:

object  myWind;

gTextOut(myWind, 10, 30, "Hello World");
@end group
@end example
@sp 1
See also:  @code{Write, Puts, Printf}
@end deffn









@deffn {Update} Update::Window
@sp 2
@example
@group
r = gUpdate(wind);

object  wind;   /*  a window object     */
object  r;      /*  wind                */
@end group
@end example
This method is used to explicitly update the display with changes the program
may have made to window @code{wind}.  WDS normally handles this need, however,
this method is available in case the programmer performs special processing
and wishes to update the entire window at one time.

This method returns the window argument passed.
@code{Update::MainWindow}.
@example
@group
@exdent Example:

object  myWind;

gUpdate(myWind);
@end group
@end example
@c @sp 1
@c See also:  @code{}
@end deffn









@deffn {Use} Use::Window
@sp 2
@example
@group
r = gUse(wind, obj);

object  wind;   /*  a window object     */
object  obj;    /*  arbitrary object    */
object  r;      /*  the object passed   */
@end group
@end example
This method is used as a general purpose mechanism to associate a
number of object types with a window.  @code{obj} may be a @code{Font,
Icon, Cursor, Menu,} or the background @code{Brush} object.  Those
objects would have normally been created via their associated classes
and then may be associated with a window via this method.

Note that the same object should not be associated with more than one
window.  The problem is that if one window is deleted, WDS would delete
all the objects associated with that window and then the other window
would reference objects which have been deleted.  The way around this is
to use the @code{Copy} method in order to make a copy of an object prior
to associating with a new window.  This way there would be two
independent objects such that if one is deleted the other would still
exist.

Note also that the @code{Application} class may be used to set application
wide defaults for these types of objects.

The value returned is the object passed.
@example
@group
@exdent Example:

object  myWind, myOtherWind, someFont;

someFont = vNew(ExternalFont, "Times New Roman", 12);
if (someFont)  @{
        gUse(myWind, someFont);
        gUse(myOtherWind, gCopy(someFont));
@}
@end group
@end example
@sp 1
See also:  @code{LoadFont, LoadIcon, LoadCursor, LoadMenu,}
        @code{TextBrush, BackBrush}
@end deffn








@deffn {VertShift} VertShift::Window
@sp 2
@example
@group
r = gVertShift(wind, rows);

object  wind;   /*  a window object     */
int     rows;   /*  rows to shift       */
object  r;      /*  the window object   */
@end group
@end example
This method is used to perform a vertical shift of the text in the
memory buffer which is being displayed in the window.  WDS keeps
a buffer which holds a programmer definable number of lines of text.
The user is then able to scroll through this text.  If the application
attempts to display more lines than this maximum (via @code{vPrintf}
for example) the system will automatically call @code{gVertShift} in
order to eliminate the top line and make room for the new line being
appended.  This method is mainly used internally.

@code{wind} is the window which is to be affected, and @code{rows} is the
number of rows to scroll.  If @code{rows} is positive the scroll moves
the text up, and negative moves the text down.  Scrolling up permanently
destroys lines at the beginning of the buffer and scrolling down
permanently destroys lines at the end of the buffer.

The scaling factor associated with @code{rows} is set by
@code{SetScalingMode::Application} and the maximum lines associated with
a window may be set with @code{SetMaxLines}.
@example
@group
@exdent Example:

object  myWind;

gVertShift(myWind, 2);
@end group
@end example
@sp 1
See also:  @code{ScrollVert, SetMaxLines}
@end deffn








@deffn {WaitCursor} WaitCursor::Window
@sp 2
@example
@group
r = gWaitCursor(wind, wait);

object  wind;   /*  a window object  */
int     wait;   /*  show/hide flag   */
object  r;      /*  wind             */
@end group
@end example
This method shows or hides the hourglass (wait) cursor for window
@code{wind} using reference counting.  Pass a non-zero value for
@code{wait} to show the wait cursor.  Pass 0 to hide it.  The
reference counting ensures that nested calls work correctly ---
the cursor is restored only when the count reaches zero.

The window object passed is returned.
@example
@group
@exdent Example:

gWaitCursor(myWind, 1);  /* show hourglass  */
/* ... long operation ... */
gWaitCursor(myWind, 0);  /* restore cursor  */
@end group
@end example
@sp 1
See also:  @code{SetCursor}
@end deffn










@deffn {Write} Write::Window
@sp 2
@example
@group
r = gWrite(wind, txt, len);

object   wind;  /*  a window object       */
char    *txt;   /*  string to be output   */
unsigned len;   /*  length of string      */
int      r;     /*  bytes written         */
@end group
@end example
This method is used to display a string of text (@code{txt}) in a sequential
fashion on window @code{wind}.  Since the @code{Window} class is a subclass
of @code{Stream}, this method is mainly provided to support the standard
interface dictated by the @code{Stream} class.  By providing this method,
the @code{Window} class automatically inherits the @code{Puts} and
@code{Printf} capability.

@code{txt} is the text to be displayed and @code{len} is the length of that
string.  Text @code{txt} will be displayed on window @code{wind} and
@code{len} will be returned.
@example
@group
@exdent Example:

object  myWind;

gWrite(myWind, "Hello World\n", 12);
@end group
@end example
@sp 1
See also:  @code{TextOut, Puts, Printf}
@end deffn



@subsection Main Window
This class, named @code{MainWindow}, is used to create and manipulate an
application's main window.  There is normally one main window associated
with each application and this is the first window created.

This class is a subclass of the @code{Window} class and therefore
inherits all of the @code{Window} class's functionality.  The methods
documented in this subsection are only those which are particular
to the @code{MainWindow} class.



@deffn {New} New::MainWindow
@sp 2
@example
@group
r = vNew(MainWindow, ttl);

char    *ttl;   /*  window title  */
object  r;      /*  new window    */
@end group
@end example
This class method is used to create the main application window.
Since this is a class method the first argument must be literally
@code{MainWindow}.  The window title (@code{ttl}) will appear at
the top of the window.

The value returned is the new window created.
@example
@group
@exdent Example:

object  myWind;

myWind = vNew(MainWindow, "My Application");
@end group
@end example
@sp 1
See also:  @code{New::ChildWindow, New::PopupWindow,}
      @code{Show, Dispose::Window}
@end deffn







@deffn {ProcessMessages} ProcessMessages::MainWindow
@sp 2
@example
@group
r = gProcessMessages(wind);

object  wind;   /*  a window object       */
int     r;      /*  final message result  */
@end group
@end example
Once the main application window is created (via @code{vNew}) and
its various attributes set, @code{gProcessMessages} is called in order
to actually create the Windows window, with the selected attributes,
display (show) it, and process the application's messages.  Processing
the application's messages is what allows the user to interact with the
application.

The value returned is the return value specified when the application is
terminated.  This can be specified via @code{QuitApplication::Application}.
@example
@group
@exdent Example:

object  myWind;

myWind = vNew(MainWindow, "Application Name");
gProcessMessages(myWind);
@end group
@end example
@sp 1
See also:  @code{New, Show::Window, ProcessMessages::MessageDispatcher}
@end deffn











@subsection Child Windows
The Child Window class, called @code{ChildWindow}, is used in the creation
of windows which are children of other windows.  As such, they can not
be positioned outside of their parent window, and they get iconized along
with their parent window.


This class is a subclass of the @code{Window} class and therefore
inherits all of the @code{Window} class's functionality.  The methods
documented in this subsection are only those which are particular
to the @code{ChildWindow} class.





@deffn {New} New::ChildWindow
@sp 2
@example
@group
r = vNew(ChildWindow, prnt, rows, cols);

object  prnt;   /*  parent window object  */
int     rows;   /*  length of window      */
int     cols;   /*  width of window       */
object  r;      /*  new window            */
@end group
@end example
This class method is used to create a child window.
Since this is a class method the first argument must be literally
@code{ChildWindow}.  @code{prnt} represents the parent window
object of which the new window will be a child.  @code{rows} and
@code{cols} determine the initial size of the window are specified
in increments dictated by @code{SetScalingMode::Application}.

The default style associated with the child window is
@code{WS_CHILD | WS_VISIBLE} and may be changed with @code{SetStyle::Window}.
The position may be set with 
@iftex
@hfil @break 
@end iftex
@code{SetPosition::Window}.

The value returned is the new window created.
@example
@group
@exdent Example:

object  myWind, mainWind;

myWind = vNew(ChildWindow, mainWind, 10, 45);
@end group
@end example
@sp 1
See also:  @code{New::MainWindow, New::PopupWindow,}
      @code{Show, Dispose::Window}
@end deffn








@subsection Popup Windows
The @code{PopupWindow} class is used to create arbitrary windows which
function independently of other windows.  That is, they may overlap or
be moved outside of other windows and iconized independently of other
windows.  If, however, a popup window is associated with a parent window
(via @code{SetParent::Window}) then it may move outside the parent but
will be iconized with it.

This class is a subclass of the @code{Window} class and therefore
inherits all of the @code{Window} class's functionality.  The methods
documented in this subsection are only those which are particular
to the @code{PopupWindow} class.






@deffn {New} New::PopupWindow
@sp 2
@example
@group
r = vNew(PopupWindow, name, rows, cols);

char    *name;  /*  window title      */
int     rows;   /*  length of window  */
int     cols;   /*  width of window   */
object  r;      /*  new window        */
@end group
@end example
This class method is used to create a popup window.  Since this is a
class method the first argument must be literally @code{PopupWindow}.
@code{name} represents the window title and will appear in the window's
top border.  @code{rows} and @code{cols} determine the initial size of
the window are specified in increments dictated by
@code{SetScalingMode::Application}.

The default style associated with the popup window is @code{WS_POPUP |
WS_VISIBLE | WS_CAPTION | WS_THICKFRAME | WS_MINIMIZEBOX |
WS_MAXIMIZEBOX | WS_SYSMENU} 
@iftex
@hfil @break 
@end iftex
and may be changed with
@code{SetStyle::Window}.  The position may be set with
@iftex
@break 
@end iftex
@code{SetPosition::Window}.

The value returned is the new window created.
@example
@group
@exdent Example:

object  myWind, mainWind;

myWind = vNew(PopupWindow, "Window Title", 10, 45);
@end group
@end example
@sp 1
See also:  @code{New::MainWindow, New::ChildWindow,}
      @code{Show, Dispose::Window}
@end deffn




